%!TEX program = xelatex
\documentclass[lang=cn,a4paper,newtx]{elegantpaper}

\title{ICU约束下基于隔离率的流行病阈值控制策略分析}
\author{ XXX \and XXX}
\institute{School of Mathematics and Statistics, Shaanxi Normal University}
% \version{0.1}
\date{}

\usepackage{amsmath,amssymb,mathrsfs}
\usepackage{booktabs}
\usepackage{graphicx}
\usepackage{float}
\usepackage{placeins}
\usepackage{array}
\usepackage{tabularx}
\usepackage{multirow}
\usepackage{makecell}
\usepackage{threeparttable}
\usepackage{siunitx}
\usepackage{xurl}
\usepackage{algorithm}
\usepackage{algpseudocode}

\graphicspath{{./}{figures/}{table/}}
\numberwithin{equation}{section}
% 推论与定理共用编号。
\theoremstyle{plain}
\newtheorem{corollaryn}[theorem]{推论}
% 需要交叉引用的注记采用编号环境。
\newtheorem{remarkn}[theorem]{注记}
\allowdisplaybreaks[3]
\renewcommand{\arraystretch}{1.12}
\emergencystretch=2em
\newcommand{\qinf}{q_{\mathrm{inf}}}
\newcommand{\Itcum}{I_{t_{\mathrm{cum}}}}
\addbibresource{references.bib}

\begin{document}

\maketitle

\begin{abstract}
\textit{（摘要待撰写。）}
\end{abstract}

\section{引言}
\textit{（引言与文献综述待撰写。）}

\section{模型与阈值控制策略}
\label{sec:model}
\subsection{包含接触追踪与隔离的扩展 SIR 型模型}
将人群分为易感者（$S$）、隔离易感者（$S_q$）、非隔离感染者（$I$）、隔离感染者（$I_q$）和康复者（$R$）五个仓室。记传播概率为 $\beta$，接触率为 $c(t)$，接触追踪比例为 $q(t)$。被追踪到的接触者中，未感染者进入 $S_q$，感染者进入 $I_q$；未被追踪到的接触者占比为 $1-q(t)$，其中感染者进入 $I$。非隔离感染者和隔离感染者分别以速率 $\gamma$ 和 $\delta_q$ 康复，并进入 $R$。

该模型为含接触追踪与隔离的扩展 SIR 型模型，下文简称 SIQR 型模型，其方程为


\begin{equation}
\begin{cases}
\dfrac{dS(t)}{dt} = -\dfrac{\beta c(t) + c(t)q(t)(1-\beta)}{N} S(t)I(t), \\[.8em] % Add some space after the first line
\dfrac{dI(t)}{dt} = \dfrac{\beta c(t)(1-q(t))}{N} S(t)I(t)  - \gamma I(t), \\[.8em]
\dfrac{dS_q(t)}{dt} = \dfrac{(1-\beta)c(t)q(t)}{N} S(t)I(t), \\[.8em]
\dfrac{dI_q(t)}{dt} = \dfrac{\beta c(t) q(t)}{N} S(t)I(t) - \delta_q I_q(t), \\[.8em]
\dfrac{dR(t)}{dt} = \gamma I(t) + \delta_q I_q(t).
\end{cases}
\label{eq:model:full}
\end{equation}

式~\eqref{eq:model:full} 中，$q(t)$ 同时影响新感染者进入 $I$ 与 $I_q$ 的比例，以及易感者进入隔离的数量。因此，即使接触率保持不变，调整隔离率仍会同时改变 $S$ 和 $I$ 的演化。相应的仓室转移结构见图~\ref{fig:model_schematic}。

\begin{figure}[htbp]
    \centering
    \includegraphics[width=0.7\textwidth]{figures/Fig1.pdf}
    \caption{含接触追踪与隔离的 SIQR 型模型。$S,S_q,I,I_q,R$ 分别表示社区易感者、隔离易感者、非隔离感染者、隔离感染者和康复者。}
    \label{fig:model_schematic}
\end{figure}
\subsection{医疗容量约束与三阶段控制策略}
本文以非隔离感染者人数的上限 $I(t)\le\eta$ 表示医疗容量约束，其中 $\eta$ 为给定的等效容量阈值，并不直接等同于实际 ICU 床位数。以感染人数上限表征医疗容量、并在达到上限后使感染人数维持在阈值水平的处理可见 Miclo 等的研究~\cite{Miclo2022ICU}。在本文模型中，接触率与隔离率均影响传播过程；为考察仅加强接触追踪隔离时的控制效果，固定 $c(t)\equiv c_0$，以 $q(t)$ 作为唯一的控制变量。

容量约束 $I(t)\le\eta$ 本身并不能唯一确定隔离率。本文选取一类以感染人数首次达到阈值为启动条件、以维持 $I(t)\equiv\eta$ 为控制目标的控制策略。记常规隔离率为 $0\le q_0<1$，常规接触率为 $c_0>0$，加强隔离的启动和结束时间分别为 $t_1$、$t_2$，阈值控制期隔离率记为 $q_c(t)$。控制函数为
\begin{equation}
c(t)\equiv c_0,\qquad
q(t)=
\begin{cases}
q_0, & t<t_1,\\
q_c(t), & t_1\le t\le t_2,\\
q_0, & t>t_2.
\end{cases}
\label{eq:model:piecewise-control}
\end{equation}

给定初值 $(S_0,I_0)$ 且 $0<I_0<\eta$，系统首先在常规水平 $(c_0,q_0)$ 下演化。若常规控制下的感染峰值不超过 $\eta$，则无需启动加强隔离；若该峰值超过 $\eta$，则在感染人数上升至阈值的首个时刻启动控制，即 $I(t_1)=\eta$，并记 $S^*=S(t_1)$。因此，$t_1$ 由常规传播过程和容量阈值共同确定，而不是预先指定的干预日期。

首次达到阈值时，常规控制下的感染人数仍在增加，继续采用 $q_0$ 将使 $I(t)$ 超过 $\eta$。本文假设隔离率可即时调整，在 $t_1$ 提高隔离率，并在此后的控制期内使进入非隔离感染仓室的流量与移出流量相等，即
\[
\frac{\beta c_0(1-q_c(t))}{N}S(t)=\gamma.
\]
由此在 $[t_1,t_2]$ 内维持 $I(t)\equiv\eta$，本文将这一阶段称为阈值控制期。阈值控制期内要求 $q_0\le q_c(t)\le1$。由于 $S(t)$ 仍随感染和隔离而减少，维持上述平衡所需的隔离率也随之变化，而不是一个预先给定的常数。

加强隔离的结束条件由恢复常规控制后感染人数是否继续增长决定。当 $q(t)=q_0$ 时，感染人数的增长率由 $\beta c_0(1-q_0)S/N-\gamma$ 的符号确定。因而，阈值控制期在易感者人数降至常规控制下的临界值时结束：
\begin{equation}
S(t_2)=S_c\doteq\frac{\gamma N}{\beta c_0(1-q_0)}.
\end{equation}
在本文不考虑易感者回流的模型中，恢复 $q_0$ 后 $S(t)$ 继续下降，因此 $t>t_2$ 时感染人数下降，不再超过容量阈值。相反，在阈值控制期内 $S>S_c$ 时恢复 $q_0$，感染人数会重新增加，因而不能满足所要求的上限。

本文将上述方案称为隔离率阈值控制。其启动和结束时刻分别由感染人数阈值与易感者临界值确定，控制持续时间是策略作用下的结果，而非独立指定的参数。以下研究这一给定策略的控制函数、持续时间、累计感染和成本；使感染人数维持在阈值水平是本文选定的控制要求，不作其为所有允许控制中最优方案的结论。

\subsection{常规阶段的首次积分与感染峰值}
判断阈值控制是否被触发，需要先确定常规控制下的感染峰值及其对应轨迹。式~\eqref{eq:model:full} 中，$S$ 和 $I$ 的方程不依赖于 $S_q$、$I_q$ 和 $R$，因此这一分析可归结为二维系统
\begin{equation}
\begin{cases}
\dot{S}(t) = - \beta_1(t) S(t) I(t) \\
\dot{I}(t) =   \beta_2(t) S(t) I(t) - \gamma I(t)
\end{cases}
\label{eq:model:reduced-si}
\end{equation}
其中
\begin{equation}
\beta_1(t)=\frac{c(t)\big[\beta+q(t)(1-\beta)\big]}{N},
\qquad
\beta_2(t)=\frac{\beta c(t)(1-q(t))}{N}.
\label{eq:model:beta12-general}
\end{equation}
其中 $\gamma$ 为非隔离感染者的移出率。由式~\eqref{eq:model:piecewise-control}，有
\begin{equation}
\beta_i(t)=
\begin{cases}
\beta_i^0, & t<t_1,\\
\beta_i^1(t), & t_1\le t\le t_2,\\
\beta_i^0, & t>t_2,
\end{cases}
\qquad i=1,2,
\end{equation}
其中常规防控阶段的系数为
\begin{equation}
\beta_1^0=\frac{c_0\big[\beta+q_0(1-\beta)\big]}{N},
\qquad
\beta_2^0=\frac{\beta c_0(1-q_0)}{N}.
\label{eq:model:beta12-baseline}
\end{equation}
加强隔离阶段的 $\beta_1^1(t)$ 与 $\beta_2^1(t)$ 由固定接触率 $c_0$ 和隔离控制函数 $q_c(t)$ 确定。

在前后两个常规控制阶段，$\beta_1$ 和 $\beta_2$ 均为常数，可通过消去时间变量求得首次积分。记
\begin{equation}
\rho_1=\frac{\gamma}{\beta_1^0},
\qquad
S_c=\frac{\gamma}{\beta_2^0}
=\frac{\gamma N}{\beta c_0(1-q_0)}.
\label{eq:model:baseline-symbols}
\end{equation}
其中 $S_c$ 是常规防控水平 $(c_0,q_0)$ 下的临界易感人数。给定初始状态 $(S_0,I_0)$，常规防控阶段的首次积分为
\begin{equation}
I(t)+\frac{\beta_2^0}{\beta_1^0} S(t)-\rho_1\ln S(t)
=
I_0+\frac{\beta_2^0}{\beta_1^0} S_0-\rho_1\ln S_0
\doteq C_0.
\label{eq:model:first-integral}
\end{equation}
在 $S_0>S_c$ 时，常规防控轨道上的感染人数可写为
\begin{equation*}
I(S)=I_0+\frac{\beta_2^0}{\beta_1^0}(S_0-S)
+\rho_1\ln\frac{S}{S_0}.
\end{equation*}
感染峰值在 $S=S_c$ 时取得，其值为
\begin{equation}
I_{\max}^{no}
=
I_0+\frac{\beta_2^0}{\beta_1^0}(S_0-S_c)
+\rho_1\ln\frac{S_c}{S_0}=C_0-\rho_1+\rho_1\ln S_c.
\label{eq:s1:Imax-no}
\end{equation}
将 $I_{\max}^{no}$ 与 $\eta$ 比较，即可判断常规控制是否已满足容量约束。以下主要讨论 $S_0>S_c$ 且 $0<I_0<\eta<I_{\max}^{no}$ 的情形，此时常规轨迹在达到峰值之前首次经过 $I=\eta$，加强隔离由此启动。

\section{隔离率阈值控制的解析结果}
\label{sec:s1}
常规阶段的首次积分确定了进入阈值控制期时的状态。阈值控制开始后，感染人数保持不变，但易感者人数和隔离率仍随时间变化；当系统达到退出条件后，感染人数又按常规控制下降。本节将这三个阶段连接起来，给出控制过程及其终止指标的解析表达。

按照式~\eqref{eq:model:piecewise-control}，整个过程中
\begin{equation*}
c(t)=c_0,\qquad t\ge 0,
\end{equation*}
隔离率在控制前后均为 $q_0$，阈值控制期内为 $q_c(t)$。

相应的二维系统为
\begin{equation}
\begin{cases}
\dot{S}(t) = - \beta_1(t) S(t) I(t) \\
\dot{I}(t) = \beta_2(t) S(t) I(t) - \gamma I(t)
\end{cases}
\label{eq:s1:dynamics}
\end{equation}
其中
\begin{equation}
\beta_i(t)=
\begin{cases}
\beta_i^0, & t<t_1,\\
\beta_i^1(t), & t_1\le t\le t_2,\\
\beta_i^0, & t>t_2,
\end{cases}
\quad i=1,2,
\end{equation}
其中 $\beta_1^0$ 和 $\beta_2^0$ 由式~\eqref{eq:model:beta12-baseline} 给出，阈值控制期的系数为
\begin{equation}
\beta_1^1(t)=\frac{c_0\big[\beta+q_c(t)(1-\beta)\big]}{N},
\quad
\beta_2^1(t)=\frac{\beta c_0(1-q_c(t))}{N}.
\label{eq:s1:beta12-control}
\end{equation}
记
\begin{equation}
\bar S
=
\frac{\gamma N(1-\beta)}{\beta c_0}.
\label{eq:s1:barS}
\end{equation}
由于
\[
\frac{\bar S}{S_c}
=
(1-\beta)(1-q_0)<1,
\]
故 $\bar S<S_c$。

\subsection{控制启动点与启动时间}
当 $I_0<\eta<I_{\max}^{no}$ 时，同一条常规轨迹在感染上升和下降阶段各经过一次 $I=\eta$。策略使用的是上升阶段的交点，因此求取 $S^*$ 时必须区分这两个交点。确定 $S^*$ 后，再由常规阶段的易感者方程得到启动时间。

\begin{theorem}
\label{thm:s1:trigger}
设 $S_0>S_c$ 且 $I_0<\eta<I_{\max}^{no}$。记
\begin{equation}
z_\eta:=
-\frac{1}{S_c}
\exp\left[-\frac{C_0-\eta}{\rho_1}\right] \quad z_\eta\in\left(-\frac{1}{e},0\right)
\label{eq:s1:lambert-argument}
\end{equation}
则常规控制下感染人数首次达到 $\eta$ 时，对应的易感者人数为
\begin{equation}
S^*=-S_cW_{-1}(z_\eta),
\qquad
S_c<S^*<S_0.
\label{eq:s1:Sstar-lambert}
\end{equation}
该启动状态唯一，相应的启动时间 $t_1>0$ 满足
\begin{equation}
t_1
=
\int_{S_0}^{S^*}
\frac{1}{
-\beta_1^0S\left[
C_0-\frac{\beta_2^0}{\beta_1^0}S+\rho_1\ln S
\right]
}\,dS.
\label{eq:s1:t1}
\end{equation}
\end{theorem}
\begin{proof}
加强控制启动前，系统满足首次积分~\eqref{eq:model:first-integral}。将 $I(t_1)=\eta$ 和 $S^*=S(t_1)$ 代入，得
\[
\frac{\beta_2^0}{\beta_1^0}S^*-\rho_1\ln S^*
=C_0-\eta.
\]
注意到
\[
\frac{\beta_2^0}{\beta_1^0\rho_1}
=
\frac{\beta_2^0}{\gamma}
=
\frac{1}{S_c}.
\]
于是上式可化为
\[
-\frac{1}{S_c}S^*
\exp\left(-\frac{1}{S_c}S^*\right)
=
-\frac{1}{S_c}
\exp\left[-\frac{C_0-\eta}{\rho_1}\right].
\]
由式~\eqref{eq:s1:lambert-argument}，上述方程可写为
\[
-\frac{S^*}{S_c}
\exp\left(-\frac{S^*}{S_c}\right)
=z_\eta.
\]
由 $I_{\max}^{no}=C_0-\rho_1+\rho_1\ln S_c$，可得 $z_{\eta}=-\frac{1}{e}\exp\left[-\frac{I_{\max}^{no}-\eta}{\rho_1}\right]$。结合 $\eta<I_{\max}^{no}$，有 $z_\eta\in(-1/e,0)$。因此，Lambert $W$ 函数存在两个实分支，且
\[
W_0(z_\eta)\in(-1,0),
\qquad
W_{-1}(z_\eta)\in(-\infty,-1).
\]
因此
\[
S^*=-S_cW_0(z_\eta)\in(0,S_c),
\]
或
\[
S^*=-S_cW_{-1}(z_\eta)>S_c.
\]
前者对应同一常规防控轨道在感染峰值后的下降阶段交点，后者对应感染上升阶段交点。
另一方面，在 $S\in(S_c,S_0)$ 上有
$\dot I=I(\beta_2^0S-\gamma)>0$。
结合 $I_0<\eta<I_{\max}^{no}$，常规防控轨道在到达峰值前恰有一次满足
$I(t)=\eta$，且该时刻满足 $S_c<S(t)<S_0$。
因此，首次达到阈值时对应的根唯一，应选择 $W_{-1}$ 分支，从而得到
式~\eqref{eq:s1:Sstar-lambert}。

确定 $S^*$ 后，可由常规阶段的易感者方程计算 $t_1$。由式~\eqref{eq:model:first-integral}，有
\[
I(S)=
C_0-\frac{\beta_2^0}{\beta_1^0}S+\rho_1\ln S,
\]
因此
\[
t_1
=
\int_{S_0}^{S^*}
\frac{1}{
-\beta_1^0S\left[
C_0-\frac{\beta_2^0}{\beta_1^0}S+\rho_1\ln S
\right]
}\,dS.
\]

\end{proof}

\begin{remark}
定理~\ref{thm:s1:trigger} 对应持续时间为正的阈值控制。若 $\eta=I_0$，则 $t_1=0$、$S^*=S_0$；若 $\eta=I_{\max}^{no}$，则 $z_\eta=-1/e$、$W_{-1}(z_\eta)=-1$、$S^*=S_c$，控制持续时间为零。若 $\eta>I_{\max}^{no}$，常规控制轨迹不会达到阈值；若 $\eta<I_0$，初始感染人数已超过阈值，不属于本文讨论的触发情形。
\end{remark}

\subsection{阈值控制期的隔离控制函数}
定理~\ref{thm:s1:trigger} 给出的启动状态满足 $S^*>S_c$。此时常规隔离率不足以阻止感染人数继续上升，必须提高 $q(t)$。恒定感染要求将这一选择具体化：进入 $I$ 的流量须恰好补偿其移出量，由此确定控制期的隔离率。

\begin{theorem}
\label{thm:s1:qc-structure}
设 $c(t)\equiv c_0$，且 $S_0>S_c$、$0<I_0<\eta<I_{\max}^{no}$。
若在阈值控制期 $t\in[t_1,t_2]$ 内要求 $I(t)\equiv\eta$，
则所需隔离率由下式唯一确定：
\begin{equation}
q(t)=q_c(S(t)),\qquad
q_c(S)=1-\frac{\gamma N}{\beta c_0S},\qquad S\in[S_c,S^*].
\label{eq:s1:q-control}
\end{equation}
结合控制前后采用的常规隔离率，完整的隔离控制函数为
\begin{equation}
q(t)=
\begin{cases}
q_0, & t<t_1,\\[0.35em]
q_c(S(t))=1-\dfrac{\gamma N}{\beta c_0S(t)}, & t_1\le t\le t_2,\\[0.7em]
q_0, & t>t_2.
\end{cases}
\label{eq:s1:q-piecewise}
\end{equation}
\end{theorem}

\begin{proof}
由恒定感染条件 $I(t)\equiv\eta$，有 $\dot I=0$。
代入系统~\eqref{eq:s1:dynamics} 的第二个方程，得
\[
\frac{\beta c_0(1-q)S}{N}\eta-\gamma\eta=0.
\]
由于 $\eta>0$，约去 $\eta$ 后解得
\[
q=1-\frac{\gamma N}{\beta c_0S}=q_c(S).
\]
该解唯一。结合三阶段控制的定义，即得式~\eqref{eq:s1:q-piecewise}。
\end{proof}

式~\eqref{eq:s1:q-control} 将所需隔离率表示为易感者人数的函数。沿阈值控制期轨迹考察这一关系，可进一步确定隔离率在启动和结束时是否连续，以及控制期内如何变化。这些性质如下。

\begin{corollaryn}
\label{cor:s1:qc-structure}
在定理~\ref{thm:s1:qc-structure} 的条件下，隔离控制函数具有以下性质：
\begin{enumerate}
\renewcommand{\labelenumi}{(\roman{enumi})}
    \item 启动时隔离率向上跳跃，即
    \[
    q(t_1^-)=q_0,\qquad
    q(t_1)=q(t_1^+)=q_{\max}
    =1-\frac{\gamma N}{\beta c_0S^*}>q_0.
    \]
    \item $q(t)$ 在 $(t_1,t_2)$ 上严格单调递减。
    \item 结束时隔离率连续，即
    \[
    q(t_2^-)=q(t_2)=q(t_2^+)=q_0.
    \]
\end{enumerate}
\end{corollaryn}

\begin{proof}
由式~\eqref{eq:s1:q-control} 和 $S_c=\gamma N/[\beta c_0(1-q_0)]$，有
\[
q_c(S_c)=q_0,\qquad q_c(S^*)=q_{\max}.
\]
在 $t=t_1$ 左侧系统仍采用常规隔离率，故 $q(t_1^-)=q_0$。
进入阈值控制期后，$q(t_1)=q(t_1^+)=q_c(S^*)=q_{\max}$，且
\[
\Delta q:=q_{\max}-q_0
=\frac{\gamma N}{\beta c_0}
\left(\frac{1}{S_c}-\frac{1}{S^*}\right)>0,
\]
其中 $S^*>S_c$，故 (i) 成立。

将 $q_c(S)$ 代入 $S$ 方程，得
\[
\dot S=-\frac{c_0\eta}{N}(S-\bar S),\qquad
\bar S=\frac{(1-\beta)\gamma N}{\beta c_0}.
\]
因此，沿阈值控制期轨迹有
\[
\frac{d}{dt}q_c(S(t))
=\frac{\gamma N}{\beta c_0}\frac{\dot S(t)}{S(t)^2}
=-\frac{\gamma\eta}{\beta}\frac{S(t)-\bar S}{S(t)^2}<0,
\qquad t_1<t<t_2,
\]
由于 $S(t)\ge S_c>\bar S$，(ii) 成立。

由 $S(t_2)=S_c$，得 $q(t_2^-)=q(t_2)=q_c(S_c)=q_0$。
而 $t>t_2$ 时隔离率恢复为 $q_0$，因此 $q(t_2^+)=q_0$，(iii) 成立。
\end{proof}

\begin{remark}
式~\eqref{eq:s1:q-control} 给出隔离控制函数的状态表达。代入下一小节的理论轨迹后，可得时间函数 $q_c(t)$；本文数值计算采用该预定函数，不作实时反馈修正。
\end{remark}

\subsection{阈值控制期的状态轨迹与控制持续时间}
上述隔离控制函数给出了隔离率与状态的关系，但确定其时间变化和持续时间还需要求出 $S(t)$。代入恒定感染条件后，易感者方程化为 $\dot S=-(c_0\eta/N)(S-\bar S)$，因此阈值控制期轨迹可以显式求得，并据此将隔离控制函数写为时间的函数。
\begin{theorem}
    在定理~\ref{thm:s1:qc-structure} 的条件下，阈值控制期 $t\in[t_1,t_2]$ 内的易感者人数为
    \begin{equation}
        S(t) = \left( S^* - \bar S \right) e^{\frac{-c_0\eta}{N}(t - t_1 )} + \bar S
        \label{eq:s1:S-analytic}
    \end{equation}
    隔离控制函数为
    \begin{equation*}
    q(t) =
    \begin{cases}
    q_0, & t < t_1, \\
    q_c(t) & t_1 \le t \le t_2, \\
    q_0, & t > t_2.       
    \end{cases}
    \end{equation*}
    其中
    \begin{equation}
        q_c(t)= 1 - \frac{1}{\left[\frac{\beta c_0 S^* }{\gamma N} + \beta- 1 \right] e^{\frac{-c_0 \eta}{N} (t-t_1)} + 1-\beta} 
    \end{equation}
\end{theorem}
\begin{proof}
    求解一阶线性方程 $\dot{S}=-\frac{c_0\eta}{N}(S-\bar S)$，得
    \[
        S(t) = C e^{\frac{-c_0 \eta}{N} t} + \bar S
    \]
    其中 $C$ 为积分常数。代入初值 $S(t_1)=S^*$，得
    \[
    S(t) = \left( S^* - \bar S \right) e^{\frac{-c_0\eta}{N}(t - t_1 )} + \bar S
    \]
    将 $S(t)$ 代入 $q_c(S(t))=1-\frac{\gamma N}{\beta c_0 S(t)}$，得到关于时间的控制函数：
    \[
    q_c(t)= 1 - \frac{1}{\left[\frac{\beta c_0 S^* }{\gamma N} + \beta- 1 \right] e^{\frac{-c_0 \eta}{N} (t-t_1)} + 1-\beta}
    \]
\end{proof}
式~\eqref{eq:s1:S-analytic} 描述了阈值控制期内易感者人数的下降过程。加强隔离并不在预定日期结束，而是持续到 $S(t)$ 达到 $S_c$。由于 $S^*>S_c>\bar S$，该轨迹在有限时间内达到退出状态，所需时间由下面的结果给出。

\begin{theorem}
    在定理~\ref{thm:s1:qc-structure} 的条件下，以 $S(t_2)=S_c$ 为结束条件，则控制结束时间为
    \begin{equation}
    t_2
    =
    t_1+
    \frac{N}{c_0 \eta}
    \ln \left[
    \frac{S^*-\bar S}
    {S_c-\bar S}
    \right]
    \label{eq:s1:t2}
    \end{equation}
    相应的控制持续时间为
        \begin{equation}
        \Delta t = \frac{N}{c_0 \eta}
        \ln \left[
        \frac{S^*-\bar S}
        {S_c-\bar S}
        \right]
    \end{equation}
\end{theorem}
\begin{proof}   
    阈值控制期内，易感者人数从 $S^*$ 下降至 $S_c$。对易感者方程积分，得
    \begin{align*}
        \Delta t=t_2-t_1=\int_{t_1}^{t_2} dt=&\int_{S^*}^{S_c} \frac{1}{\dot{S}} \, dS
        =\frac{N}{c_0\eta}\int_{S_c}^{S^*} \frac{1}{S-\bar S}\,dS
        =\frac{N}{c_0 \eta}
        \ln \left[
        \frac{S^*-\bar S}
        {S_c-\bar S}
        \right].
    \end{align*}
    因此，
    \begin{equation*}
    t_2
    =
    t_1+
    \frac{N}{c_0 \eta}
    \ln \left[
    \frac{S^*-\bar S}
    {S_c-\bar S}
    \right]=\int_{S_0}^{S^*}
    \frac{1}{
    -\beta_1^0S\left[
    C_0-\frac{\beta_2^0}{\beta_1^0}S+\rho_1\ln S
    \right]
    }\,dS + \frac{N}{c_0 \eta}
    \ln \left[
    \frac{\frac{\beta c_0}{\gamma N}S^*+ \beta-1}
    {\frac{q_0}{1-q_0}+ \beta}
    \right]
    \end{equation*}
\end{proof}

\begin{remark}
    将结束条件 $S(t_2)=S_c=\frac{\gamma N}{\beta c_0(1-q_0)}$ 代入式~\eqref{eq:s1:S-analytic}，亦得
    \begin{equation*}
     t_2= t_1+\frac{N}{c_0 \eta}
     \ln \left[
     \frac{S^*-\bar S}
     {S_c-\bar S}
     \right]
    \end{equation*}
    以及控制持续时间
    \begin{equation*}
    \Delta t = \frac{N}{c_0 \eta}
     \ln \left[
     \frac{S^*-\bar S}
     {S_c-\bar S}
     \right]
    \end{equation*}
\end{remark}
\begin{remark}
    在 $S(t)=C e^{\frac{-c_0 \eta}{N}t}+\bar S$ 中代入结束条件 $S(t_2)=S_c$，可写为
    \begin{equation*}
    S(t) = \left( S_c - \bar S \right) e^{\frac{c_0\eta}{N}(t_2 - t )} + \bar S
    \end{equation*}
    再代入 $q_c(S(t))=1-\frac{\gamma N}{\beta c_0S(t)}$，隔离控制函数可写为
    \begin{align*}
    q_c(t)=1 - \frac{1}{(\tfrac{q_0}{1-q_0}+\beta)e^{\frac{c_0\eta}{N}(t_2 - t )}+1-\beta}
    \end{align*}
\end{remark}

\subsection{动态清零时间}
控制结束并不等于感染过程结束：在 $t_2$ 时仍有 $I(t_2)=\eta$。恢复常规隔离率后，$S(t)<S_c$，感染人数随之下降；当 $\eta>1$ 时，还需要一段时间才能达到社区动态清零。因此，控制持续时间与疫情持续时间应分别计算。

沿用 Zhou 等~\cite{Zhou2024SPCI} 对隔离区外动态清零的定义，本文以非隔离感染者人数在下降阶段达到 $1$ 作为判据，记相应时刻为 $t_{end}$，即 $I(t_{end})=1$ 且 $I'(t_{end})<0$。退出后的常规轨迹从 $(S_c,\eta)$ 出发，利用同一首次积分即可确定该时刻及其对应的易感者人数。
\begin{theorem}\label{thm:s1:tend}
    在上述阈值控制下，设动态清零时刻位于控制结束之后。记 $C_2=\eta+(\beta_2^0/\beta_1^0)S_c-\rho_1\ln S_c$，则 $t_{end}$ 由下式给出
    \[
    t_{end}=t_2+\int_{S_c}^{S(t_{end})}\frac{1}
    {S\left[\beta_2^0 S-\beta_1^0 \rho_1 \ln(S)-
    \beta_1^0 C_2\right]}\, dS .
    \]
    其中
    \begin{equation}
    S(t_{end})=
    -\frac{\beta_1^0 \rho_1}{\beta_2^0}\,
    \operatorname{LambertW}_0
    \left(
    -\frac{\beta_2^0}{\beta_1^0 \rho_1}
    \exp\left(
    \frac{1-C_2}{\rho_1}
    \right)
    \right).
    \label{eq:s1:Send}
    \end{equation}
\end{theorem}
\begin{proof}
    由于在阈值控制中，当 $t\ge t_2$ 时，系统恢复到常规防控水平
    \[
    c(t)=c_0,\qquad q(t)=q_0.
    \]
    因此，$t\in[t_2,+\infty)$ 上仍满足常规阶段的首次积分。不过此时首次积分常数由
    $t=t_2$ 时刻的状态决定。记
    \[
    C_2=\eta+\frac{\beta_2^0}{\beta_1^0}S(t_2)-\rho_1\ln S(t_2)=\eta+\frac{\beta_2^0}{\beta_1^0}S_c-\rho_1\ln S_c.
    \]
    于是
    \[
    I(t)
    =
    -\frac{\beta_2^0}{\beta_1^0} S(t)
    +
    \rho_1 \ln S(t)
    +
    C_2,
    \qquad
    t \in [t_2,+\infty).
    \]
    由 $I(t_{end})=1$ 和 $I'(t_{end})<0$ 得到
    \[
    1
    +
    \frac{\beta_2^0}{\beta_1^0} S(t_{end})
    -
    \rho_1 \ln(S(t_{end}))
    =
    C_2 .
    \]
    同时满足 $S(t_{end})<\frac{\gamma}{\beta_2^0}$。利用 Lambert $W$ 函数求解，得到
    \[
    S(t_{end})=
    -\frac{\beta_1^0 \rho_1}{\beta_2^0}\,
    \operatorname{LambertW}_0
    \left(
    -\frac{\beta_2^0}{\beta_1^0 \rho_1}
    \exp\left(
    \frac{1-C_2}{\rho_1}
    \right)
    \right).
    \]
    于是
    \[
    S'(t)=S\left[\beta_2^0 S-
    \beta_1^0 \rho_1 \ln(S)
    -\beta_1^0 C_2
    \right],
    \qquad
    t \in [t_2,+\infty).
    \]
    从时间 $t_2$ 积分到 $t_{end}$，有
    \[
    \int_{S_c}^{S(t_{end})}
    \frac{1}
    {
    S\left[
    \beta_2^0 S
    -
    \beta_1^0 \rho_1 \ln(S)
    -
    \beta_1^0 C_2
    \right]
    }
    \, dS
    =
    \int_{t_2}^{t_{end}} dt .
    \]

    因此
    \[
    t_{end}=t_2+\int_{S_c}^{S(t_{end})}\frac{1}
    {S\left[\beta_2^0 S-\beta_1^0 \rho_1 \ln(S)-
    \beta_1^0 C_2\right]}\, dS .
    \]
\end{proof}


\subsection{总累计感染}
峰值约束只限制同时存在的非隔离感染者人数，不能由此确定整个过程的累计感染规模。阈值控制期虽有 $I(t)\equiv\eta$，仍不断出现新感染者；此外，直接进入 $I_q$ 的感染者也应计入总数。为与动态清零时间采用同一评价终点，将从初始时刻到动态清零时刻的总累计感染定义为
\begin{equation}
\Itcum
=
\int_0^{t_{\rm end}}\frac{\beta c(t)}{N}S(t)I(t)\,dt.
\end{equation}
同时，进入社区感染仓室 $I$ 与进入隔离感染仓室 $I_q$ 的累计感染分别为
\begin{equation}
I_{\rm cum}
=
\int_0^{t_{\rm end}}
\frac{\beta c(t)(1-q(t))}{N}S(t)I(t)\,dt,
\end{equation}
和
\begin{equation}
I_{q_{\rm cum}}
=
\int_0^{t_{\rm end}}
\frac{\beta c(t)q(t)}{N}S(t)I(t)\,dt.
\end{equation}
于是
\[
\Itcum=I_{\rm cum}+I_{q_{\rm cum}}.
\]

前后两个常规阶段具有相同的参数，感染流量均可用易感者人数的变化表示；阈值控制阶段则使用已求得的 $S(t)$ 和 $q_c(t)$。分别计算三段流量，可将累计感染写成启动、退出和清零时刻的易感者人数的函数。

\begin{theorem}\label{thm:s1:Itcum}
    设系统在 $t_1$ 启动阈值控制，在 $t_2$ 结束控制，并在 $t_{\rm end}>t_2$ 达到动态清零。记 $S_{\rm end}=S(t_{\rm end})$，则从初始时刻到 $t_{\rm end}$ 的总累计感染为
    \[
    \Itcum
    =
    \frac{\beta}{\beta+q_0(1-\beta)}(S_0-S^*)
    +
    \beta\left[
    (S^*-S_c)+\bar S\ln\frac{S^*-\bar S}{S_c-\bar S}
    \right]
    +
    \frac{\beta}{\beta+q_0(1-\beta)}(S_c-S_{\rm end}).
    \]
\end{theorem}
\begin{proof}
    记总新发感染率为
    \[
    F(t)=\frac{\beta c(t)}{N}S(t)I(t).
    \]
    在第一段常规阶段 $0\le t\le t_1$，有 $c(t)=c_0,\ q(t)=q_0$，故
    \[
    \dot S
    =
    -\frac{c_0[\beta+q_0(1-\beta)]}{N}SI.
    \]
    因而
    \[
    F(t)
    =
    \frac{\beta}{\beta+q_0(1-\beta)}(-\dot S).
    \]
    从 $0$ 到 $t_1$ 积分，并利用 $S(t_1)=S^*$，得到
    \[
    \int_0^{t_1}F(t)\,dt
    =
    \frac{\beta}{\beta+q_0(1-\beta)}(S_0-S^*).
    \]

    在阈值控制阶段 $t_1\le t\le t_2$，有 $I(t)\equiv\eta$，且
    \[
    q_c(t)=1-\frac{\gamma N}{\beta c_0S(t)}.
    \]
    由 $S$ 方程可得
    \[
    \dot S
    =
    -\frac{c_0\eta}{N}(S-\bar S),
    \qquad
    \bar S=\frac{(1-\beta)\gamma N}{\beta c_0}.
    \]
    因此
    \[
    dt
    =
    -\frac{N}{c_0\eta}\frac{dS}{S-\bar S}.
    \]
    又因为此时
    \[
    F(t)=\frac{\beta c_0}{N}S(t)\eta,
    \]
    所以
    \[
    \int_{t_1}^{t_2}F(t)\,dt
    =
    \beta\int_{S_c}^{S^*}\frac{S}{S-\bar S}\,dS
    =
    \beta\left[
    (S^*-S_c)+\bar S\ln\frac{S^*-\bar S}{S_c-\bar S}
    \right].
    \]

    在解除控制后的常规阶段 $t_2\le t\le t_{\rm end}$，仍有
    $c(t)=c_0,\ q(t)=q_0$，因此与第一段相同，
    \[
    F(t)
    =
    \frac{\beta}{\beta+q_0(1-\beta)}(-\dot S).
    \]
    由 $S(t_2)=S_c$ 与 $S(t_{\rm end})=S_{\rm end}$ 得
    \[
    \int_{t_2}^{t_{\rm end}}F(t)\,dt
    =
    \frac{\beta}{\beta+q_0(1-\beta)}(S_c-S_{\rm end}).
    \]

    三段相加即得结论。
\end{proof}

\section{成本与参数敏感性}
\label{sec:cost}
第~\ref{sec:s1}~节确定了给定参数下的三阶段控制过程。不同容量阈值和常规接触水平会改变启动状态、所需隔离率及控制持续时间，因而即使都能满足各自的感染上限，控制投入也可能不同。本节以累计控制强度评价这一投入，并分析有关指标的参数依赖。

以下分析均在相应结论规定的参数范围内进行。对某一参数求偏导时，除特别说明外，固定人口规模、初始状态及其余独立参数。

\subsection{成本函数}

为与其他控制方案采用统一指标进行比较，本文保留接触率和隔离率两个分项归一化成本
\begin{equation}
J_c=\int_{t_1}^{t_2} \frac{\bigl(c_0-c(t)\bigr)_+}{c_0}\,dt,
\end{equation}
\begin{equation}
J_q=\int_{t_1}^{t_2} \frac{\bigl(q(t)-q_0\bigr)_+}{1-q_0}\,dt.
\end{equation}
其中 $(x)_+=\max\{x,0\}$。进一步定义二次加权综合成本
\begin{equation}
J
=
\int_{t_1}^{t_2}
\left[
\left(\frac{\bigl(c_0-c(t)\bigr)_+}{c_0}\right)^2
+
2\left(\frac{\bigl(q(t)-q_0\bigr)_+}{1-q_0}\right)^2
\right]dt.
\label{eq:cost:J}
\end{equation}
式~\eqref{eq:cost:J} 的被积函数为归一化控制强度，不显含人口规模 $N$。

隔离项的权重取为 $2$，表示在该成本设定下加强隔离的单位强度成本较高。本文阈值策略保持 $c(t)\equiv c_0$，故 $J_c=0$，额外成本仅来自隔离率的提高。
利用控制期内
\[
\dot S=-\frac{c_0\eta}{N}(S-\bar S),
\]
可将阈值控制中的 $J_q$ 写为
\begin{equation}
J_q
=
\frac{N}{\eta(1-q_0)}
\int_{S_c}^{S^*}
\frac{q_c(S)-q_0}{c_0(S-\bar S)}\,dS.
\end{equation}
当 $\eta$ 增大时，上限 $S^*$ 下降，同时前面的 $1/\eta$ 下降，故额外累计控制强度随
$\eta$ 增大而下降。对于 $c_0$，积分上限、终点阈值、被积函数以及时间尺度同时改变，
这些变化的作用方向不同，需结合后续数值结果判断。


\subsection{启动点的参数敏感性}
控制持续时间、最大隔离率和成本的表达式都含有 $S^*$，而 $S^*$ 又由常规轨迹与感染阈值的交点确定。因此，分析这些指标时需要同时计入启动状态的变化，不能把 $S^*$ 视为独立于参数的常量。下面先给出其参数导数的符号，作为后续分析的依据。

\begin{lemma}
\label{lem:s1:Sstar-sensitivity}
设 $S_0>S_c$ 且 $I_0<\eta<I_{\max}^{no}$。固定其余参数时，
启动点 $S^*$ 满足
\[
\frac{\partial S^*}{\partial \eta}<0,
\qquad
\frac{\partial S^*}{\partial c_0}>0,
\qquad
\frac{\partial S^*}{\partial \beta}>0,
\qquad
\frac{\partial S^*}{\partial q_0}<0.
\]
若固定 $N,S_0,I_0$ 并以 $\theta=\eta/N$ 为参数，则
\[
\frac{\partial S^*}{\partial \theta}=N\frac{\partial S^*}{\partial \eta}<0.
\]
\end{lemma}

\begin{proof}
由常规控制阶段的首次积分，启动点满足
\begin{equation}
\begin{aligned}
F(S^*;\eta,c_0,\beta,q_0)
:={}&
I_0+
\frac{\beta(1-q_0)}
{\beta+q_0(1-\beta)}
(S_0-S^*) \\
&+
\rho_1\ln\frac{S^*}{S_0}
-\eta
=0.
\end{aligned}
\label{eq:s1:F-Sstar-param}
\end{equation}
由于 $\rho_1/S_c=\beta(1-q_0)/[\beta+q_0(1-\beta)]$，
\[
\frac{\partial F}{\partial S^*}
=
\rho_1
\left(
\frac1{S^*}-\frac1{S_c}
\right)<0,
\]
由于启动时感染人数处于上升阶段，有 $S^*>S_c$，故上述偏导数非零，隐函数定理适用。

首先，$\frac{\partial F}{\partial \eta}=-1$，故
\begin{equation}
\frac{\partial S^*}{\partial \eta}
=
\frac{1}{
\rho_1
\left(
\dfrac1{S^*}-\dfrac1{S_c}
\right)
}
<0.
\label{eq:s1:Sstar-eta-sign}
\end{equation}
其次，$\frac{\partial \rho_1}{\partial c_0}=-\rho_1/c_0$，从而
\[
\frac{\partial F}{\partial c_0}
=
-\frac{\rho_1}{c_0}\ln\frac{S^*}{S_0}>0,
\]
并有
\begin{equation}
\frac{\partial S^*}{\partial c_0}
=
\frac{\ln(S^*/S_0)}
{c_0
\left(
\dfrac1{S^*}-\dfrac1{S_c}
\right)
}
>0.
\label{eq:s1:Sstar-c0-sign}
\end{equation}

对 $\beta$，固定 $S^*$ 时
\[
\frac{\partial}{\partial\beta}
\frac{\beta(1-q_0)}
{\beta+q_0(1-\beta)}
=
\frac{q_0(1-q_0)}
{\left[\beta+q_0(1-\beta)\right]^2},
\]
\[
\frac{\partial \rho_1}{\partial \beta}
=
-\frac{\rho_1(1-q_0)}
{\beta+q_0(1-\beta)}.
\]
因此
\begin{equation}
\frac{\partial S^*}{\partial \beta}
=
\frac{
\displaystyle
\frac{q_0(1-q_0)}
{\left[\beta+q_0(1-\beta)\right]^2}
(S_0-S^*)
-
\frac{\rho_1(1-q_0)}
{\beta+q_0(1-\beta)}
\ln\frac{S^*}{S_0}
}{
\displaystyle
\rho_1
\left(
\frac1{S_c}-\frac1{S^*}
\right)
}
>0.
\label{eq:s1:Sstar-beta-sign}
\end{equation}
上式分母为正。由 $S_0>S^*$、$\ln(S^*/S_0)<0$，分子第一项非负、第二项严格为正，故 $\frac{\partial S^*}{\partial \beta}>0$。

最后，对 $q_0$ 有
\[
\frac{\partial}{\partial q_0}
\frac{\beta(1-q_0)}
{\beta+q_0(1-\beta)}
=
-\frac{\beta}
{\left[\beta+q_0(1-\beta)\right]^2},
\qquad
\frac{\partial \rho_1}{\partial q_0}
=
-\frac{\rho_1(1-\beta)}
{\beta+q_0(1-\beta)}.
\]
利用
\[
\frac{\rho_1(1-\beta)}
{\beta+q_0(1-\beta)}
=
\frac{\beta\bar S}
{\left[\beta+q_0(1-\beta)\right]^2},
\]
可得
\begin{equation}
\frac{\partial S^*}{\partial q_0}
=
-\frac{
\beta
\left[
(S_0-S^*)
-\bar S\ln\frac{S_0}{S^*}
\right]
}{
\left[\beta+q_0(1-\beta)\right]^2
\rho_1
\left(
\dfrac1{S_c}-\dfrac1{S^*}
\right)
}
<0.
\label{eq:s1:Sstar-q0-sign}
\end{equation}
事实上，由 $\bar S<S^*<S_0$ 及 $\ln x<x-1$（$x>1$）知
\[
\bar S\ln\frac{S_0}{S^*}
<
S^*\ln\frac{S_0}{S^*}
<
S_0-S^*,
\]
故式~\eqref{eq:s1:Sstar-q0-sign} 方括号内严格为正。
当固定 $N$ 且 $\eta=N\theta$ 时，链式法则给出
$\frac{\partial S^*}{\partial \theta}=N\frac{\partial S^*}{\partial \eta}<0$。
\end{proof}

\subsection{启动时间与最大隔离率的参数敏感性}
启动点的变化反映加强隔离时的易感者规模，启动时间还取决于到达该状态的速度；两者的变化方向不能直接等同。最大隔离率则由启动状态和传播参数共同决定。以下分别考察常规接触率与医疗容量阈值
\[
c_0,\qquad \eta .
\]
的影响，固定其余参数，并限定在常规轨迹于感染上升阶段达到阈值的范围内。

为便于书写，记
\begin{equation}
a=\frac{\beta_2^0}{\beta_1^0}
=
\frac{\beta(1-q_0)}{\beta+q_0(1-\beta)},
\qquad
\rho_1=\frac{\gamma}{\beta_1^0}
=
\frac{\gamma N}{c_0[\beta+q_0(1-\beta)]},
\end{equation}
以及
\begin{equation}
S_c=\frac{\gamma}{\beta_2^0}
=
\frac{\gamma N}{\beta c_0(1-q_0)}.
\end{equation}
注意到 $a$ 与 $c_0$ 无关，而 $\rho_1$ 与 $S_c$ 均与 $c_0$ 成反比。

改变 $c_0$ 还可能使常规感染峰值降至阈值以下，此时加强隔离不再被触发。因此，在比较启动时间和隔离强度之前，先确定阈值控制被触发的参数范围。固定 $\eta$ 及 $N,S_0,I_0,\beta,\gamma,q_0$，记常规控制下的感染峰值为 $I_{\max}^{no}(c_0)$，定义临界接触率
\begin{equation}
c_{0,\min}(\eta)
=
\inf\left\{
c_0>0:
I_{\max}^{no}(c_0)\geq \eta
\right\}.
\label{eq:s1:c0-min-def}
\end{equation}
由式~\eqref{eq:s1:Imax-no}，并利用
$\beta_2^0/\beta_1^0=a$、$\rho_1=aS_c$，可将常规控制下的峰值写为
\begin{equation}
I_{\max}^{no}(c_0)
=
I_0+a(S_0-S_c)+aS_c\ln\frac{S_c}{S_0}.
\label{eq:s1:Imax-no-c0}
\end{equation}
当 $S_0>S_c$、感染人数在初始阶段增加时，
\begin{equation}
\frac{d I_{\max}^{no}}{dc_0}
=
-\frac{aS_c}{c_0}\ln\frac{S_c}{S_0}>0.
\label{eq:s1:Imax-no-c0-derivative}
\end{equation}
若给定阈值满足
$I_0<\eta<I_0+aS_0$，式~\eqref{eq:s1:c0-min-def} 所定义的临界值在上述参数范围内存在且唯一，并满足
\begin{equation}
I_{\max}^{no}\!\left(c_{0,\min}(\eta)\right)=\eta.
\label{eq:s1:c0-min-equality}
\end{equation}
后文固定 $\eta$ 时，将 $c_{0,\min}(\eta)$ 简记为 $c_{0,\min}$。于是
\begin{equation}
\begin{aligned}
c_0<c_{0,\min}
&\quad\Longrightarrow\quad I_{\max}^{no}(c_0)<\eta,\\
c_0=c_{0,\min}
&\quad\Longrightarrow\quad I_{\max}^{no}(c_0)=\eta,\\
c_0>c_{0,\min}
&\quad\Longrightarrow\quad I_{\max}^{no}(c_0)>\eta.
\end{aligned}
\label{eq:s1:c0-min-regimes}
\end{equation}
当 $c_0<c_{0,\min}$ 时，无需加强隔离；当 $c_0=c_{0,\min}$ 时，轨迹在感染峰值处与 $I=\eta$ 相切，$S^*=S_c$ 且 $\Delta t=0$；当 $c_0>c_{0,\min}$ 时，控制持续时间为正。
相应地，
\begin{equation}
c_0\downarrow c_{0,\min}^{+}
\quad\Longrightarrow\quad
S^*\to S_c,
\qquad
\Delta t\to0.
\label{eq:s1:c0-min-limit}
\end{equation}
在 $c_0>c_{0,\min}$ 时，$\Delta t$ 的变化方向由推论~\ref{cor:s1:duration-sensitivity} 给出的导数判据确定。

\begin{proposition}
\label{prop:s1:sensitivity}
设参数满足以下触发条件：
\[
c_0>c_{0,\min}(\eta),
\qquad
I_0<\eta<I_{\max}^{no}(c_0),
\qquad
S_0>S^*>S_c>\bar S.
\]
固定其余参数，启动时间和最大隔离率满足
\[
\frac{\partial t_1}{\partial \eta}>0,
\qquad
\frac{\partial t_1}{\partial c_0}<0,
\]
以及
\[
\frac{\partial q_{\max}}{\partial \eta}<0,
\qquad
\frac{\partial q_{\max}}{\partial c_0}>0.
\]
\end{proposition}

\begin{proof}
由引理~\ref{lem:s1:Sstar-sensitivity}，有
$\frac{\partial S^*}{\partial \eta}<0$ 与 $\frac{\partial S^*}{\partial c_0}>0$。启动时间可写为
\begin{equation}
t_1
=
\frac{N}{c_0\bigl[\beta+q_0(1-\beta)\bigr]}
\int_{S^*}^{S_0}
\frac{1}{S I(S;c_0)}\,dS,
\qquad
I(S;c_0)
=
I_0-a(S-S_0)+\rho_1\ln\frac{S}{S_0}.
\end{equation}
由于积分下端满足 $I(S^*;c_0)=\eta$，对 $\eta$ 求导得到
\begin{equation}
\frac{\partial t_1}{\partial \eta}
=
-\frac{N}{c_0\bigl[\beta+q_0(1-\beta)\bigr]}
\frac{1}{S^*\eta}
\frac{\partial S^*}{\partial \eta}>0.
\end{equation}
对 $c_0$ 求导，有
\[
\frac{\partial I(S;c_0)}{\partial c_0}
=
-\frac{\rho_1}{c_0}\ln\frac{S}{S_0}>0,
\qquad S\in(S^*,S_0),
\]
从而
\begin{align}
\frac{\partial t_1}{\partial c_0}
=&
-\frac{N}{c_0^2\bigl[\beta+q_0(1-\beta)\bigr]}
\int_{S^*}^{S_0}\frac{1}{S I(S;c_0)}\,dS
-\frac{N}{c_0\bigl[\beta+q_0(1-\beta)\bigr]}
\frac{1}{S^*\eta}
\frac{\partial S^*}{\partial c_0} \nonumber\\
&-\frac{N}{c_0\bigl[\beta+q_0(1-\beta)\bigr]}
\int_{S^*}^{S_0}
\frac{\frac{\partial I(S;c_0)}{\partial c_0}}{S I^2(S;c_0)}\,dS
<0.
\end{align}


由隔离控制函数的表达式
\[
q_c(S)=1-\frac{\gamma N}{\beta c_0S},
\qquad
\frac{dq_c}{dS}
=
\frac{\gamma N}{\beta c_0S^2}>0.
\]
控制期内 $S(t)$ 单调下降，因此 $q_c(t)$ 在控制期内单调下降，最大值出现在启动时刻 $t_1$：
\[
q_{\max}
=
q_c(t_1)
=
1-\frac{\gamma N}{\beta c_0S^*}.
\]
于是
\[
\frac{\partial q_{\max}}{\partial \eta}
=
\frac{\gamma N}{\beta c_0(S^*)^2}\frac{\partial S^*}{\partial \eta}<0,
\]
以及
\[
\frac{\partial q_{\max}}{\partial c_0}
=
\frac{\gamma N}{\beta c_0^2(S^*)^2}
\left(S^*+c_0\frac{\partial S^*}{\partial c_0}\right)>0.
\]
\end{proof}

\subsection{控制时长的参数敏感性}
命题~\ref{prop:s1:sensitivity} 表明，增大 $c_0$ 会使控制更早启动，并提高启动时所需的隔离率，但这并不直接决定阈值控制需持续多久。式~\eqref{eq:s1:t2} 中，参数既改变起止状态，也改变时间尺度，控制时长的导数需要同时考虑这些作用。结合启动点的参数导数，可得以下结果。

\begin{corollaryn}
\label{cor:s1:duration-sensitivity}
设 $S_0>S_c$ 且 $I_0<\eta<I_{\max}^{no}$。固定其余参数时，有
\[
\frac{\partial \Delta t}{\partial \eta}<0,
\qquad
\frac{\partial \Delta t}{\partial \beta}>0,
\qquad
\frac{\partial \Delta t}{\partial q_0}<0.
\]
若固定 $N,S_0,I_0$，则
$\frac{\partial \Delta t}{\partial \theta}=N\frac{\partial \Delta t}{\partial \eta}<0$。
另一方面，$\frac{\partial \Delta t}{\partial c_0}$ 没有固定符号。令
\[
\Theta
=
\frac{S^*-S_c}{S_c}
\left[
\frac{S^*-\bar S}{S^*}
\ln\frac{S^*-\bar S}{S_c-\bar S}
-1
\right],
\]
则
\[
\frac{\partial\Delta t}{\partial c_0}
\gtrless0
\quad\Longleftrightarrow\quad
\ln\frac{S_0}{S^*}
\gtrless\Theta,
\]
等号情形亦对应成立。特别地，若 $\Theta\le0$，则
$\partial\Delta t/\partial c_0>0$。
\end{corollaryn}

\begin{proof}
由
\[
\Delta t
=
\frac{N}{c_0\eta}
\ln\frac{S^*-\bar S}{S_c-\bar S}
\]
及 $\frac{\partial S^*}{\partial \eta}<0$，有
\[
\frac{\partial\Delta t}{\partial\eta}
=
-\frac{N}{c_0\eta^2}
\ln\frac{S^*-\bar S}{S_c-\bar S}
+
\frac{N}{c_0\eta}
\frac{\frac{\partial S^*}{\partial \eta}}{S^*-\bar S}
<0.
\]
固定 $N$ 时 $\eta=N\theta$，故
$\frac{\partial \Delta t}{\partial \theta}=N\frac{\partial \Delta t}{\partial \eta}<0$；等价地，
\[
\frac{\partial\Delta t}{\partial\theta}
=
\frac1{c_0}
\left[
\frac1{\theta}
\frac{\frac{\partial S^*}{\partial \theta}}{S^*-\bar S}
-
\frac1{\theta^2}
\ln\frac{S^*-\bar S}{S_c-\bar S}
\right]<0.
\]

对 $\beta$，
\[
\frac{\partial}{\partial\beta}
\ln\frac{S^*-\bar S}{S_c-\bar S}
=
\frac{\frac{\partial S^*}{\partial \beta}-\frac{\partial \bar S}{\partial \beta}}{S^*-\bar S}
-
\frac{\frac{\partial S_c}{\partial \beta}-\frac{\partial \bar S}{\partial \beta}}{S_c-\bar S}.
\]
其中
\[
\frac{\partial \bar S}{\partial \beta}
=
-\frac{\bar S}{\beta(1-\beta)}<0,
\qquad
\frac{\partial S_c}{\partial \beta}-\frac{\partial \bar S}{\partial \beta}
=
-\frac{q_0S_c}{\beta}\le0.
\]
第一项严格为正，第二项非负，故 $\frac{\partial \Delta t}{\partial \beta}>0$。

由于 $\bar S$ 不含 $q_0$，
\[
\frac{\partial}{\partial q_0}
\ln\frac{S^*-\bar S}{S_c-\bar S}
=
\frac{\frac{\partial S^*}{\partial q_0}}{S^*-\bar S}
-
\frac{\frac{\partial S_c}{\partial q_0}}{S_c-\bar S},
\qquad
\frac{\partial S_c}{\partial q_0}
=
\frac{S_c}{1-q_0}>0.
\]
结合 $\frac{\partial S^*}{\partial q_0}<0$，得到 $\frac{\partial \Delta t}{\partial q_0}<0$。

最后，由
$\frac{\partial \bar S}{\partial c_0}=-\bar S/c_0$、
$\frac{\partial S_c}{\partial c_0}=-S_c/c_0$ 及
式~\eqref{eq:s1:Sstar-c0-sign}，可得
\[
\frac{\partial\Delta t}{\partial c_0}
=
\frac{N}{c_0^2\eta}
\left[
\frac{S^*}{S^*-\bar S}
\left(
1+
\frac{S_c}{S^*-S_c}
\ln\frac{S_0}{S^*}
\right)
-
\ln\frac{S^*-\bar S}{S_c-\bar S}
\right].
\]
整理方括号中的不等式，得
$\ln(S_0/S^*)\gtrless\Theta$。由于
$\ln(S_0/S^*)>0$，当 $\Theta\le0$ 时导数必为正。
\end{proof}

在上述触发范围内，提高 $\eta$ 对应更晚的启动时刻、更低的最大隔离率和更短的阈值控制期，其代价是允许更高的非隔离感染人数。相比之下，增大 $c_0$ 虽然提高最大隔离率，却不一定延长阈值控制期。这些结果区分了控制强度与持续时间的变化；累计感染和动态清零时间还包含退出后的感染过程，不能仅由控制持续时间推断。

\section{隔离控制函数的拐点分析}
\label{sec:s1:shape}
前面的结果给出了控制持续时间和最大隔离率，并证明 $q_c(t)$ 在控制期内严格递减。然而，单调性并不说明隔离率的下降速度是否一直增大或减小。考察内部拐点可以进一步确定下降速度是否在阈值控制期内达到最大值，以及这一时刻如何随参数变化。

\subsection{拐点的存在条件}
\begin{theorem}\label{thm:s1:inflection}
    设 $S_0>S_c$ 且 $I_0<\eta<I_{\max}^{no}$。隔离控制函数 $q_c(t)$ 在 $(t_1,t_2)$ 内存在唯一拐点，当且仅当
    \[
    S_c<2\bar S<S^*.
    \]
    等价地，
    \[
    \frac{1}{2}
    <(1-\beta)(1-q_0)
    <-\frac{1}{2}W_{-1}(z_\eta).
    \]
    当内部拐点存在时，其对应的易感人数和发生时间分别为：
    \[
    S(t_{\mathrm{inf}})=2\bar S,\quad
    t_{\mathrm{inf}}
    =
    t_1+\frac{N}{c_0\eta}\ln\frac{S^*-\bar S}{\bar S},
    \]
    拐点处的隔离率为 $\qinf:=q_c(2\bar S)=1-\dfrac{1}{2(1-\beta)}$，仅依赖于 $\beta$。

    此外，
    \[
    q_c''(t)<0,
    \qquad t_1<t<t_{\mathrm{inf}},
    \]
    而
    \[
    q_c''(t)>0,
    \qquad t_{\mathrm{inf}}<t<t_2.
    \]
\end{theorem}
\begin{proof}
    对推论~\ref{cor:s1:qc-structure} 中沿阈值控制期轨迹的一阶导数继续求导，得
    \[
    q_c''(t)
    =
    \frac{\frac{\gamma c_0\eta^2}{\beta N}(S-\bar S)(2\bar S-S)}{S^3}.
    \]
    由于 \(\frac{\gamma c_0\eta^2}{\beta N}>0,\ S>\bar S\)，所以 $q_c''(t)$ 的符号只由 $2\bar S-S(t)$ 决定：
    \[
    S(t)>2\bar S
    \quad\Rightarrow\quad
    q_c''(t)<0,
    \quad S(t)<2\bar S
    \quad\Rightarrow\quad
    q_c''(t)>0.
    \]
    由于 $S(t)$ 在阈值控制期由 $S^*$ 严格递减至 $S_c$，方程 $S(t)=2\bar S$ 在 $(t_1,t_2)$ 内至多有一个解，且存在当且仅当
    \[
    S_c<2\bar S<S^*.
    \]
    由此有
    \[
    \frac{2\bar{S}}{S_c}=2(1-\beta)(1-q_0)
    \]
    故 $2\bar S>S_c$ 等价于 $(1-\beta)(1-q_0)>\frac{1}{2}$。
    
    同时，由式~\eqref{eq:s1:Sstar-lambert}，
    \[
    \frac{S^*}{2\bar S}
    =
    -\frac{W_{-1}(z_\eta)}
    {2(1-\beta)(1-q_0)}.
    \]
    所以 $2\bar S<S^*$ 等价于
    \[
    (1-\beta)(1-q_0)<-\frac{1}{2}W_{-1}(z_\eta).
    \]

    综上，内部拐点 $t_{\mathrm{inf}}$ 存在，当且仅当
    \[
    \frac{1}{2}
    <(1-\beta)(1-q_0)
    <-\frac{1}{2}W_{-1}(z_\eta),
    \]
    此时 $S(t_{\mathrm{inf}})=2\bar S$。
    
    由 $\dot{S}=-\frac{c_0\eta}{N}(S-\bar S)$，求解并代入初值条件 $S(t_1)=S^*$，可得
    \[
    S(t)=\bar S+(S^*-\bar S)
    e^{-\frac{c_0\eta}{N}(t-t_1)}
    \]

    又因为\(S(t_{\mathrm{inf}})=2\bar S\)，得到
    \[
    t_{\mathrm{inf}}
    =
    t_1+
    \frac{N}{c_0\eta}
    \ln\frac{S^*-\bar S}{\bar S}.
    \]
\end{proof}
由于 $q_c(S)=1-\gamma N/(\beta c_0 S)$ 关于 $S$ 严格递增，定理~\ref{thm:s1:inflection} 中的拐点存在条件 $S_c<2\bar S<S^*$ 可等价地用隔离率表示（其中 $\qinf=q_c(2\bar S)$）：
\begin{equation}
q_0<\qinf<q_{\max}.
\label{eq:s1:inflection-exist}
\end{equation}
其中 $\qinf>q_0$ 等价于 $(1-\beta)(1-q_0)>1/2$，仅取决于 $\beta,q_0$；$\qinf<q_{\max}$ 则通过启动点 $S^*$ 依赖于 $c_0,\eta$ 和初值。


\subsection{拐点的相对位置与边界情形}
\label{sec:s1:inflection-position}
定理~\ref{thm:s1:inflection} 给出了拐点的绝对时刻，但不同参数下阈值控制的起止时间也会改变。若要比较隔离率下降最快的时刻在各自控制区间内的位置，需要将拐点时间与控制持续时间联系起来。

当 $S_c<2\bar S<S^*$ 时，由阈值控制期轨迹
\[
S(t)=\bar S+(S^*-\bar S)e^{-c_0\eta(t-t_1)/N},
\]
可得拐点前、后两段时长分别为
\begin{equation}
t_{\mathrm{inf}}-t_1
=
\frac{N}{c_0\eta}
\ln\frac{S^*-\bar S}{\bar S},
\label{eq:s1:seg-high}
\end{equation}
\begin{equation}
t_2-t_{\mathrm{inf}}
=
\frac{N}{c_0\eta}
\ln\frac{\bar S}{S_c-\bar S}.
\label{eq:s1:seg-low}
\end{equation}
两式相加即为控制时长 $\Delta t$（式~\eqref{eq:s1:t2}）。

定义拐点在控制区间内的归一化相对位置
\begin{equation}
\lambda
:=
\frac{t_{\mathrm{inf}}-t_1}{t_2-t_1}
=
\frac{
\ln\dfrac{S^*-\bar S}{\bar S}
}{
\ln\dfrac{S^*-\bar S}{S_c-\bar S}
}
\in(0,1).
\label{eq:s1:lambda}
\end{equation}
$\lambda$ 接近 $0$ 或 $1$，分别表示拐点相对靠近控制起点或终点。
拐点处 $\lvert q_c'(t)\rvert$ 最大，因此 $\lambda$ 给出隔离率下降最快时刻的相对位置，但不直接反映绝对时刻 $t_{\mathrm{inf}}$ 的变化。

上述相对位置只在内部拐点存在时使用。参数变化使拐点到达控制区间端点后，其内部存在条件不再成立。若 $2\bar S\le S_c$，则 $q_c''(t)<0$ 对所有 $t\in(t_1,t_2)$ 成立；等号 $2\bar S=S_c$ 对应拐点到达控制终点 $t_2$。
若 $2\bar S\ge S^*$，则 $q_c''(t)>0$ 对所有 $t\in(t_1,t_2)$ 成立；等号 $2\bar S=S^*$ 对应拐点到达控制起点 $t_1$。
在两种等号情形中，二阶导数仅在相应端点取零，控制区间内部均不存在拐点。

\subsection{拐点相对位置的参数敏感性}
拐点处的隔离率 $\qinf$ 仅由 $\beta$ 决定，但拐点的相对位置还取决于阈值控制的起止状态。因此，拐点处隔离率保持不变，并不意味着 $\lambda$ 不变。结合第~\ref{sec:cost}~节对 $S^*$ 的分析，可以确定各参数对 $\lambda$ 的影响。

\begin{proposition}
\label{prop:s1:lambda-sensitivity}
设 $S_0>S_c$、$I_0<\eta<I_{\max}^{no}$ 且 $S_c<2\bar S<S^*$。固定其余参数时，有
\[
\frac{\partial \lambda}{\partial \eta}<0,
\qquad
\frac{\partial \lambda}{\partial c_0}>0,
\qquad
\frac{\partial \lambda}{\partial \beta}>0.
\]
关于常规控制隔离率 $q_0$，偏导 $\frac{\partial \lambda}{\partial q_0}$ 一般没有固定符号；
其驻点满足
\begin{equation}
\frac{
\ln\dfrac{S^*-\bar S}{\bar S}
}{
(1-q_0)\bigl[\beta+q_0(1-\beta)\bigr]
}
=
-
\frac{
\ln\dfrac{\bar S}{S_c-\bar S}
}{
S^*-\bar S
}
\frac{\partial S^*}{\partial q_0}.
\label{eq:s1:lambda-q0-stationary}
\end{equation}
若固定 $N,S_0,I_0$，则
$\frac{\partial \lambda}{\partial \theta}=N\frac{\partial \lambda}{\partial \eta}<0$。
\end{proposition}

\begin{proof}
记
\[
A:=
\ln\frac{S^*-\bar S}{\bar S},
\qquad
B:=
\ln\frac{\bar S}{S_c-\bar S}.
\]
内部拐点条件保证 $A,B>0$。对所考察的任一参数 $p$，有
\[
\lambda=\frac{A}{A+B},
\qquad
\frac{\partial \lambda}{\partial p}
=
\frac{B\frac{\partial A}{\partial p}-A\frac{\partial B}{\partial p}}{(A+B)^2}.
\]
由引理~\ref{lem:s1:Sstar-sensitivity}，
\[
\frac{\partial A}{\partial \eta}
=
\frac{\frac{\partial S^*}{\partial \eta}}{S^*-\bar S}<0,
\qquad
\frac{\partial B}{\partial \eta}=0,
\]
故 $\frac{\partial \lambda}{\partial \eta}=B\frac{\partial A}{\partial \eta}/(A+B)^2<0$。

由于 $\frac{\partial \bar S}{\partial c_0}=-\bar S/c_0$，直接整理得
\[
\frac{\partial A}{\partial c_0}
=
\frac{\frac{\partial S^*}{\partial c_0}+S^*/c_0}{S^*-\bar S}>0,
\qquad
\frac{\partial B}{\partial c_0}=0,
\]
从而 $\frac{\partial \lambda}{\partial c_0}>0$。

又有
\[
\frac{\partial \bar S}{\partial \beta}
=
-\frac{\bar S}{\beta(1-\beta)},
\]
\[
\frac{\partial A}{\partial \beta}
=
\frac{
\frac{\partial S^*}{\partial \beta}+\dfrac{S^*}{\beta(1-\beta)}
}{
S^*-\bar S
}>0,
\qquad
\frac{\partial B}{\partial \beta}
=
-\frac{1}{
(1-\beta)\bigl[\beta+q_0(1-\beta)\bigr]
}<0.
\]
因此
\[
\frac{\partial \lambda}{\partial \beta}
=
\frac{B\frac{\partial A}{\partial \beta}-A\frac{\partial B}{\partial \beta}}{(A+B)^2}>0.
\]

最后，$\bar S$ 不含 $q_0$，故
\[
\frac{\partial A}{\partial q_0}
=
\frac{\frac{\partial S^*}{\partial q_0}}{S^*-\bar S}<0,
\qquad
\frac{\partial B}{\partial q_0}
=
-\frac{1}{
(1-q_0)\bigl[\beta+q_0(1-\beta)\bigr]
}<0.
\]
于是
\[
\frac{\partial \lambda}{\partial q_0}
=
\frac{1}{(A+B)^2}
\left[
\frac{B \frac{\partial S^*}{\partial q_0}}{S^*-\bar S}
+
\frac{A}{
(1-q_0)\bigl[\beta+q_0(1-\beta)\bigr]
}
\right].
\]
方括号内两项符号相反。令其和为零，并代入 $A,B$ 的表达式，即得
式~\eqref{eq:s1:lambda-q0-stationary}。固定 $N$ 时，
$\eta=N\theta$ 给出 $\frac{\partial \lambda}{\partial \theta}=N\frac{\partial \lambda}{\partial \eta}<0$。
\end{proof}


在内部拐点存在的范围内，提高 $\eta$ 使隔离率下降最快的时刻在阈值控制期内相对前移，提高 $\beta$ 或 $c_0$ 则使其相对后移。$q_0$ 的影响由相反方向的两项共同决定，不能据此断言每条参数曲线都非单调。启动时易感者人数、拐点相对位置及控制持续时间的参数敏感性汇总于表~\ref{tab:s1:lambda-sensitivity} 和表~\ref{tab:s1:inflection-roles}。两表均列出关于 $\eta$ 的结果；固定 $N$ 时，改用 $\theta=\eta/N$ 不改变相应偏导数的符号。

\begin{table}[htbp]
  \centering
  \small
  \setlength{\tabcolsep}{3.5pt}
  \renewcommand{\arraystretch}{1.25}
  \begin{threeparttable}
    \caption{启动时易感者人数与拐点相对位置的参数敏感性}
    \label{tab:s1:lambda-sensitivity}
    \begin{tabularx}{\textwidth}{@{}>{\raggedright\arraybackslash}p{0.29\textwidth}*{4}{>{\centering\arraybackslash}p{0.10\textwidth}}>{\raggedright\arraybackslash}X@{}}
      \toprule
      \multirow{2}{*}{变量} & \multicolumn{4}{c}{偏导数符号} & \multirow{2}{*}{依据}\\
      \cmidrule(lr){2-5}
       & $\beta$ & $q_0$ & $c_0$ & $\eta$ & \\
      \midrule
      启动时易感者人数 $S^*$ & $>0$ & $<0$ & $>0$ & $<0$ & 引理~\ref{lem:s1:Sstar-sensitivity}\\
      拐点相对位置 $\lambda$ & $>0$ & 不定 & $>0$ & $<0$ & 命题~\ref{prop:s1:lambda-sensitivity}\\
      \bottomrule
    \end{tabularx}
    \begin{tablenotes}[flushleft]
      \footnotesize
      \item 注：表中各项为分析量关于列示参数的偏导数符号，求导时固定 $N,S_0,I_0$ 及其余模型参数。关于 $S^*$ 的结果要求满足控制触发条件；关于 $\lambda$ 的结果还要求 $S_c<2\bar S<S^*$。不定表示偏导数的符号不能在上述参数范围内统一确定。
    \end{tablenotes}
  \end{threeparttable}
\end{table}


\begin{table}[htbp]
  \centering
  \small
  \setlength{\tabcolsep}{3.5pt}
  \renewcommand{\arraystretch}{1.25}
  \begin{threeparttable}
    \caption{拐点相关量与控制持续时间的参数敏感性}
    \label{tab:s1:inflection-roles}
    \begin{tabularx}{\textwidth}{@{}>{\raggedright\arraybackslash}p{0.29\textwidth}*{4}{>{\centering\arraybackslash}p{0.10\textwidth}}>{\raggedright\arraybackslash}X@{}}
      \toprule
      \multirow{2}{*}{变量} & \multicolumn{4}{c}{偏导数符号} & \multirow{2}{*}{依据}\\
      \cmidrule(lr){2-5}
       & $\beta$ & $q_0$ & $c_0$ & $\eta$ & \\
      \midrule
      拐点处隔离率 $\qinf$ & $<0$ & $=0$ & $=0$ & $=0$ & 定理~\ref{thm:s1:inflection}\\
      $\qinf-q_0$ & $<0$ & $<0$ & $=0$ & $=0$ & 式~\eqref{eq:s1:inflection-exist}\\
      $S^*-2\bar S$ & $>0$ & $<0$ & $>0$ & $<0$ & 引理~\ref{lem:s1:Sstar-sensitivity}、式~\eqref{eq:s1:barS}\\
      控制持续时间 $\Delta t$ & $>0$ & $<0$ & 不定 & $<0$ & 推论~\ref{cor:s1:duration-sensitivity}\\
      \bottomrule
    \end{tabularx}
    \begin{tablenotes}[flushleft]
      \footnotesize
      \item 注：表中各项为分析量关于列示参数的偏导数符号，求导时固定 $N,S_0,I_0$ 及其余模型参数。结果在控制触发条件下成立。$\qinf-q_0$ 和 $S^*-2\bar S$ 同时为正时存在内部拐点；$=0$ 表示该量不依赖相应参数。$\partial\Delta t/\partial c_0$ 的符号由推论~\ref{cor:s1:duration-sensitivity} 中的判据确定。
    \end{tablenotes}
  \end{threeparttable}
\end{table}


\section{数值验证与参数分析}
\label{sec:num}
\subsection{基准参数与解析解验证}
本节数值示例取
\[
N=763,\quad S_0=762,\quad I_0=1,\quad
\beta=0.155,\quad c_0=10,\quad q_0=0.01526,\quad \gamma=0.3504.
\]
常规控制下的再生数为
\[
R_0=\frac{\beta c_0(1-q_0)}{\gamma}\approx 4.3560.
\]

取医疗容量阈值为总人数的 $5\%$，即
$\eta=0.05N=38.15$。
表~\ref{tab:results} 和图~\ref{fig:scenario1:single-sim} 给出阈值控制与常规控制的比较。采用预定时间控制函数 $q_c(t)$ 后，$I(t)$ 在 $[t_1,t_2]$ 内维持于 $\eta$，最大数值偏差约为 $1.07\times10^{-7}$。常规控制下的感染峰值约为 $300.3726$，超过容量阈值。
\begin{figure}[ht]
    \centering
    \includegraphics[width=0.7\textwidth]{figures/optimal_control_with_quarantine_panels.pdf}
    \caption{基准参数下阈值控制与常规控制的时间轨迹。阈值控制在 $[t_1,t_2]$ 内维持 $I(t)=\eta$，隔离率由 $q_{\max}$ 逐渐下降至 $q_0$。}
    \label{fig:scenario1:single-sim}
\end{figure}

\begin{table}[htbp]
  \centering
  \small
  \sisetup{group-digits=false}
  \setlength{\tabcolsep}{3.5pt}
  \renewcommand{\arraystretch}{1.25}
  \begin{threeparttable}
    \caption{阈值控制数值模拟中的参数、理论结果与数值验证}
    \label{tab:results}
    \begin{tabularx}{\textwidth}{@{}>{\raggedright\arraybackslash}X>{\centering\arraybackslash}p{0.16\textwidth}>{\raggedright\arraybackslash}X>{\centering\arraybackslash}p{0.16\textwidth}@{}}
\toprule
变量（符号） & 数值 & 变量（符号） & 数值 \\
\midrule
\multicolumn{4}{@{}l}{模型参数} \\
\addlinespace[2pt]
移出率 $\gamma$ & 0.3504 & 基本再生数 $R_0$ & 4.3560 \\
传播概率 $\beta$ & 0.1550 & 常规接触率 $c_0$ & 10.0000 \\
总人数 $N$ & 763.0000 & 常规隔离率 $q_0$ & 0.0153 \\
医疗容量阈值 $\eta$ & 38.1500 & 临界易感人数 $S_c$ & 175.1602 \\
初始易感人数 $S_0$ & 762.0000 & 初始感染人数 $I_0$ & 1.0000 \\
\midrule
\multicolumn{4}{@{}l}{理论结果} \\
\addlinespace[2pt]
启动点 $S^*$ & 708.3470 & 启动时间 $t_1$ & 3.1754 \\
解除时间 $t_2$ & 9.0780 & 控制时长 $\Delta t$ & 5.9026 \\
最大隔离率 $q_c(t_1)$ & 0.7565 & 终点隔离强度 $q_c(t_2)=q_0$ & 0.0153 \\
\midrule
\multicolumn{4}{@{}l}{数值验证} \\
\addlinespace[2pt]
受控情形感染峰值 & 38.1500 & 阈值维持误差 & $1.07\times10^{-7}$ \\
常规控制感染峰值 & 300.3726 & 常规控制峰值时间 & 6.7065 \\
\bottomrule
    \end{tabularx}
    \begin{tablenotes}[flushleft]
      \footnotesize
      \item 注：时间单位为 d。感染峰值为相应策略下的 $\max I(t)$；阈值维持误差为 $\max_{[t_1,t_2]}|I-\eta|$。
    \end{tablenotes}
  \end{threeparttable}
\end{table}
\subsection{拐点的数值分析}

图~\ref{fig:s1:lambda-sensitivity} 给出 $\lambda$ 随 $\eta/N,c_0,\beta,q_0$ 的变化。
在内部拐点存在的范围内，$\lambda$ 随 $\eta$ 增大而减小，随 $c_0$ 和 $\beta$ 增大而增大，与命题~\ref{prop:s1:lambda-sensitivity} 一致。
固定 $N=763,S_0=762,I_0=1,\gamma=0.3504,c_0=10,\eta/N=0.05$，在 $q_0\in[0,0.39]$ 内得到以下局部极大值位置：
\[
\begin{array}{c|ccc}
\beta & 0.10 & 0.12 & 0.155\\
\hline
q_0^\dagger
&0.0736525&0.1969606&0.3890712
\end{array}
\]
这三个位置两侧的 $\frac{\partial \lambda}{\partial q_0}$ 均由正变负，且极大值位置随 $\beta$ 增大而增大。这是所考察参数范围内的数值结果，不涉及一般参数下驻点的唯一性。

\begin{figure}[htbp]
    \centering
    \includegraphics[width=\textwidth]{scenario1_inflection_lambda_sensitivity.pdf}
    \caption{拐点相对位置 $\lambda$ 随 $\eta/N,c_0,\beta,q_0$ 的变化。基准参数同本节设定；$q_0$ 面板另比较 $\beta=0.10,0.12,0.155$。空心圆表示数值局部极大值，浅灰区域表示不满足触发条件，深灰区域表示已触发控制但无内部拐点。$q_0$ 面板的灰色区域对应 $\beta=0.155$。}
    \label{fig:s1:lambda-sensitivity}
\end{figure}


图~\ref{fig:s1:inflection-landscape} 给出拐点的存在范围及其前后两段时长。改变 $\beta$ 时，两端的拐点边界均可能出现；改变 $c_0$ 时，内部拐点在 $S^*=2\bar S$ 处出现。对于图示的 $q_0$ 和 $\eta/N$ 范围，只要满足控制触发条件，内部拐点就存在。由于 $\qinf$ 仅依赖于 $\beta$，其值在其余三个参数的比较中保持不变。相应的隔离率时间轨迹见图~\ref{fig:s1:inflection-scan}。

\begin{figure}[htbp]
    \centering
    \includegraphics[width=0.8\textwidth]{scenario1_inflection_landscape.pdf}
    \caption{四个参数对拐点存在性和控制持续时间的影响，基准参数同本节设定。左列为 $q_{\max},\qinf,q_0$，满足 $q_0<\qinf<q_{\max}$ 时存在内部拐点；右列分别为拐点前后的时长 $t_{\rm inf}-t_1$ 和 $t_2-t_{\rm inf}$，虚线表示总时长 $\Delta t$。}
    \label{fig:s1:inflection-landscape}
\end{figure}
\begin{figure}[htbp]
    \centering
    \includegraphics[width=0.8\textwidth]{scenario1_inflection_scan_t.pdf}
    \caption{不同 $q_0,\beta,c_0,\eta/N$ 下的隔离率时间轨迹。其余参数同图~\ref{fig:s1:inflection-landscape}。圆点为内部拐点，短竖线为控制结束时刻 $t_2$。}
    \label{fig:s1:inflection-scan}
\end{figure}

附录~\ref{app:c0_beta} 给出了 $(c_0,\beta)$ 平面上的拐点存在范围。固定 $\theta=0.002$、$q_0=0.3230$ 时，西安参数 $(c_0,\beta)=(12.8872,0.1498)$ 位于该范围内。

\subsection{接触率与容量阈值的影响}
固定基准参数中的其余取值，分别考察 $c_0$ 和 $\eta$ 的影响。不同 $c_0$ 对应不同的常规接触水平，每条轨迹内仍保持 $c(t)\equiv c_0$。当 $\eta=0.05N$ 时，控制触发边界为
\[
c_{0,\min}(0.05N)\approx3.3270.
\]
参数范围取为
\[
c_0=2.3,2.4,\ldots,14,
\qquad
\frac{\eta}{N}=0.2\%,0.22\%,\ldots,5\%,
\]
仅在满足触发条件的参数范围内计算控制指标。表~\ref{tab:scenario1_c0_summary} 固定 $\eta/N=5\%$，表~\ref{tab:scenario1_eta_summary} 固定 $c_0=10$；相应的时间轨迹与参数变化见图~\ref{fig:scenario1_u_time_c0}--\ref{fig:scenario1_heatmaps_c0_eta}。

\begin{table}[htbp]
  \centering
  \small
  \sisetup{group-digits=false}
  \setlength{\tabcolsep}{3.5pt}
  \renewcommand{\arraystretch}{1.25}
  \begin{threeparttable}
    \caption{固定 $\eta/N=5\%$ 时不同 $c_0$ 下的关键控制指标}
    \label{tab:scenario1_c0_summary}
    \begin{tabular*}{\textwidth}{@{\extracolsep{\fill}}l*{6}{S[table-format=3.4]}@{}}
    \toprule
    \multirow{2}{*}{变量} & \multicolumn{6}{c}{常规接触率 $c_0$} \\
    \cmidrule(lr){2-7}
     & {3.347} & {4} & {5} & {8} & {10} & {12} \\
    \midrule
    $S^*$ & 559.1496 & 654.6273 & 682.5742 & 703.7125 & 708.3470 & 711.0304 \\
    $S_c$ & 523.3306 & 437.9004 & 350.3203 & 218.9502 & 175.1602 & 145.9668 \\
    \addlinespace[3pt]
    $t_1$ & 29.9176 & 15.3697 & 9.3039 & 4.3073 & 3.1754 & 2.5150 \\
    $\Delta t$ & 2.0432 & 6.8658 & 7.5778 & 6.6308 & 5.9026 & 5.3007 \\
    $t_{\rm end}$ & 71.4690 & 58.7277 & 49.9657 & 38.0095 & 33.7732 & 30.7687 \\
    \addlinespace[3pt]
    $q_{\max}$ & 0.0783 & 0.3413 & 0.4946 & 0.6936 & 0.7565 & 0.7978 \\
    $J$ & 0.0053 & 0.4374 & 1.0561 & 2.0006 & 2.2236 & 2.3110 \\
    $I_{t_{\rm cum}}$ & 382.5525 & 360.7318 & 342.5629 & 303.5829 & 285.7244 & 271.9295 \\
    \bottomrule
    \end{tabular*}
    \begin{tablenotes}[flushleft]
      \footnotesize
      \item 注：第一列参数取 $c_0=c_{0,\min}+0.02\approx3.3470$，其余列为主文代表值。
    \end{tablenotes}
  \end{threeparttable}
\end{table}

\begin{table}[htbp]
  \centering
  \small
  \sisetup{group-digits=false}
  \setlength{\tabcolsep}{3.5pt}
  \renewcommand{\arraystretch}{1.25}
  \begin{threeparttable}
    \caption{固定 $c_0=10$ 时不同医疗容量阈值下的关键控制指标}
    \label{tab:scenario1_eta_summary}
    \begin{tabular*}{\textwidth}{@{\extracolsep{\fill}}l*{4}{S[table-format=3.4]}@{}}
    \toprule
    \multirow{2}{*}{变量} & \multicolumn{4}{c}{医疗容量阈值 $\eta$} \\
    \cmidrule(lr){2-5}
     & {1.526} & {4.578} & {7.63} & {15.26} \\
    \midrule
    $\eta/N$ & 0.002 & 0.006 & 0.01 & 0.02 \\
    $S^*$ & 761.2486 & 756.8844 & 752.5126 & 741.5490 \\
    $S_c$ & 175.1602 & 175.1602 & 175.1602 & 175.1602 \\
    \addlinespace[3pt]
    $t_1$ & 0.3602 & 1.3001 & 1.7406 & 2.3459 \\
    $\Delta t$ & 152.0574 & 50.5672 & 30.2685 & 15.0431 \\
    $t_{\rm end}$ & 180.5821 & 86.5579 & 64.9343 & 46.8029 \\
    \addlinespace[3pt]
    $q_{\max}$ & 0.7734 & 0.7721 & 0.7708 & 0.7674 \\
    $J$ & 60.8005 & 20.1265 & 11.9912 & 5.8888 \\
    $I_{t_{\rm cum}}$ & 173.0314 & 195.3349 & 208.9166 & 233.7810 \\
    \bottomrule
    \end{tabular*}
    \begin{tablenotes}[flushleft]
      \footnotesize
      \item 注：各列对应 $\eta/N=0.2\%,0.6\%,1\%,2\%$。
    \end{tablenotes}
  \end{threeparttable}
\end{table}

\begin{figure}[htbp]
    \centering
    \includegraphics[width=0.7\textwidth]{figures/scenario1_u_time_c0.pdf}
    \caption{固定 $\eta/N=5\%$，$c_0=4,5,8,10,12$ 时的隔离率轨迹。三角形与短竖线分别表示 $t_1,t_2$，圆点为内部拐点。$c_0=4$ 时无内部拐点。}
    \label{fig:scenario1_u_time_c0}
\end{figure}

\begin{figure}[htbp]
    \centering
    \includegraphics[width=0.7\textwidth]{figures/scenario1_u_time_eta.pdf}
    \caption{固定 $c_0=10$，$\eta/N=0.2\%,0.6\%,1\%,2\%$ 时的隔离率轨迹。三角形与短竖线分别表示 $t_1,t_2$，圆点为内部拐点。各拐点对应相同的隔离率 $\qinf$。}
    \label{fig:scenario1_u_time_eta}
\end{figure}

\begin{figure}[htbp]
    \centering
    \includegraphics[width=0.8\textwidth]{figures/c0_sensitivity_selected_eta.pdf}
    \caption{不同阈值比例下各变量随 $c_0$ 的变化。 $\eta/N=0.2\%,0.6\%,1\%,2\%$，红色圆点表示$\Delta t$极大值。}
    \label{fig:scenario1_summary_c0}
\end{figure}

\begin{figure}[htbp]
    \centering
    \includegraphics[width=0.8\textwidth]{figures/eta_sensitivity_selected_c0.pdf}
    \caption{不同常规接触率下各变量随$\eta/N$ 的变化，$c_0=5,8,10,12$。}
    \label{fig:scenario1_summary_eta}
\end{figure}

\begin{figure}[htbp]
    \centering
    \includegraphics[width=0.8\textwidth]{figures/scenario1_heatmaps_c0_eta.pdf}
    \caption{$(\eta/N,c_0)$ 平面上的启动时间、控制持续时间、清零时间和累计感染。实线为触发边界 $c_{0,\min}(\eta)$，白色区域无需加强隔离，虚线为各指标的等值线。}
    \label{fig:scenario1_heatmaps_c0_eta}
\end{figure}

%\clearpage

图~\ref{fig:scenario1_u_time_c0} 中，$c_0=4$ 时 $q_{\max}=0.3413<\qinf\approx0.4083$，不存在内部拐点；$c_0=5,8,10,12$ 时均存在唯一内部拐点。图~\ref{fig:scenario1_u_time_eta} 中，各拐点对应的隔离率相同，但启动时间和控制持续时间随阈值改变。

由表~\ref{tab:scenario1_c0_summary} 和图~\ref{fig:scenario1_summary_c0}，增大 $c_0$ 会提高启动点 $S^*$ 和最大隔离率 $q_{\max}$，并使启动时间提前。控制时长在所考察范围内先增加后减小。例如，$\eta/N=5\%$ 时，
\[
\Delta t(3.3470)=2.0432,\qquad
\Delta t(5)=7.5778,\qquad
\Delta t(12)=5.3007\ {\rm d}.
\]
相应的数值极大值位于 $c_0\approx5.1$，约为 $7.5788$ 天。对于 $\eta/N=0.2\%,0.6\%,1\%,2\%$，极大值位置分别约为 $c_0=4.3,4.4,4.4,4.6$，对应时长为 $211.5435,69.8417,41.5169,20.2813$ 天。

表~\ref{tab:scenario1_eta_summary} 和图~\ref{fig:scenario1_summary_eta} 表明，增大 $\eta$ 会推迟控制启动、降低 $q_{\max}$，并缩短阈值控制期。当 $c_0=10$、$\eta/N$ 从 $0.2\%$ 增至 $2\%$ 时，
\[
\Delta t:152.0574\to15.0431,\qquad
t_{\rm end}:180.5821\to46.8029,
\]
\[
J:60.8005\to5.8888,\qquad
\Itcum:173.0314\to233.7810.
\]
在这一参数范围内，提高阈值缩短了阈值控制期和清零时间，降低了控制成本，但增加了累计感染。图~\ref{fig:scenario1_heatmaps_c0_eta} 给出这些指标在 $(\eta/N,c_0)$ 平面上的变化。触发边界 $c_{0,\min}(\eta)$ 随 $\eta$ 增大而升高；其下方的常规感染峰值低于阈值，无需启动加强隔离。

\section{西安疫情数据应用}
\label{sec:xian}

本节采用 He、Tang 和 Xiao~\cite{He2023TDINN} 给出的西安疫情 TDINN 控制函数和日报数据，将隔离率阈值控制用于西安参数下的数值计算，并与 TDINN 控制及常规控制比较。TDINN 控制同时降低接触率、提高隔离率；阈值控制则保持接触率不变，仅通过提高隔离率使社区感染者人数 $I(t)$ 不超过 $\eta$。记社区感染者 $I(t)$、新增社区感染 $I_{\rm new}(t)$、新增隔离感染 $I_{q_{\rm new}}(t)$、累计社区感染 $I_{\rm cum}(t)$、累计隔离感染 $I_{q_{\rm cum}}(t)$，以及总累计感染 $\Itcum(t)=I_{\rm cum}(t)+I_{q_{\rm cum}}(t)$。

\subsection{真实数据与 TDINN 控制函数}
病例数据包含社区发现病例数 $I_{\rm new}(t)$ 和隔离发现病例数 $I_{q_{\rm new}}(t)$，覆盖 2021 年 12 月 9 日至 2022 年 1 月 17 日共 40 天，其中社区发现 583 例、隔离发现 1467 例，合计 2050 例（图~\ref{fig:xian_observed_daily}）。TDINN 控制曲线取自文献~\cite{He2023TDINN} 表 3 中西安的拟合结果 $c_2(t),q_2(t)$：
\begin{equation}
c_{\rm TDINN}(t)=(12.8872-3.4625)e^{-(0.0463t)^2}+3.4625,\qquad
q_{\rm TDINN}(t)=(0.3230-0.9844)e^{-(0.0452t)^2}+0.9844,
\end{equation}
因此常规初始控制水平为 $c_0=12.8872,\ q_0=0.3230$。

\begin{figure}[htbp]
    \centering
    \includegraphics[width=0.78\textwidth]{xian_observed_daily_cases.pdf}
    \caption{西安疫情的每日新增病例。$I_{\rm new}$ 和 $I_{q_{\rm new}}$ 分别为社区发现与隔离发现病例数。}
    \label{fig:xian_observed_daily}
\end{figure}

\subsection{模型与初值拟合}
在系统~\eqref{eq:model:full} 中，引入累计变量 $I_{\rm cum}(t),I_{q_{\rm cum}}(t)$，满足
\begin{equation}
\dot I_{\rm cum}(t)=\frac{\beta c(t)(1-q(t))S(t)I(t)}{N},\qquad
\dot I_{q_{\rm cum}}(t)=\frac{\beta c(t)q(t)S(t)I(t)}{N}.
\end{equation}
总人口取西安市 2021 年末常住人口 $N=13{,}163{,}000$；传播参数固定为文献~\cite{He2023TDINN} 表 1 的估计值 $\beta=0.1498,\ \gamma=0.2953,\ \delta_q=0.3531$。由于原文未给出拟合初值，本文设 $R(0)=0$，从而 $N=S_0+I_0$，只用最小二乘估计 $I_0$ 并令 $S_0=N-I_0$；拟合准则为每日新增 $I_{\rm new},I_{q_{\rm new}}$ 与累计 $I_{\rm cum},I_{q_{\rm cum}}$ 四项均方误差之和。拟合结果见表~\ref{tab:xian_initial_fit}。

拟合得到的 $I_0$ 是给定时间原点和控制函数下的模型初值，不对应整数病例数。若改取 $I_0=1$，模型计算的 $40$ 天累计病例约为 $1.08\times10^6$，远高于观测值，见表~\ref{tab:xian_i0_sensitivity}。

\begin{table}[htbp]
  \centering
  \small
  \sisetup{group-digits=false}
  \setlength{\tabcolsep}{3.5pt}
  \renewcommand{\arraystretch}{1.25}
  \begin{threeparttable}
    \caption{固定参数后的初始条件估计}
    \label{tab:xian_initial_fit}
    \begin{tabularx}{\textwidth}{@{}>{\raggedright\arraybackslash}p{0.34\textwidth}>{\centering\arraybackslash}X S[table-format=8.6]@{}}
    \toprule
    名称 & 符号或定义 & {取值} \\
    \midrule
    总人口 & $N$ & 13163000 \\
    初始易感人数 & $S_0$ & 13162999.9990 \\
    初始感染人数 & $I_0$ & 0.001007 \\
    初始康复人数 & $R(0)=N-S_0-I_0$ & 0.0000 \\
    \bottomrule
    \end{tabularx}
    \begin{tablenotes}[flushleft]
      \footnotesize
      \item 注：$I_0$ 为模型中的有效初始感染人数；初值拟合的残差准则为均方误差（MSE）。
    \end{tablenotes}
  \end{threeparttable}
\end{table}

\begin{table}[htbp]
  \centering
  \small
  \sisetup{group-digits=false}
  \setlength{\tabcolsep}{3.5pt}
  \renewcommand{\arraystretch}{1.25}
  \begin{threeparttable}
    \caption{不同初值设定下的模型输出}
    \label{tab:xian_i0_sensitivity}
    \begin{tabular*}{\textwidth}{@{\extracolsep{\fill}}lS[table-format=1.6]S[table-format=7.0]*{2}{S[table-format=5.2]}@{}}
    \toprule
    \multirow{2}{*}{初值设定} & {\multirow{2}{*}{$I_0$}} & {\multirow{2}{*}{累计病例}} & \multicolumn{2}{c}{新增峰值} \\
    \cmidrule(lr){4-5}
     & & & {社区} & {隔离区} \\
    \midrule
    真实数据 & {---} & 2050 & 60.00 & 141.00 \\
    最小二乘拟合 & 0.001007 & 2099 & 53.88 & 115.93 \\
    强行取 $I_0=1$ & 1.000000 & 1083554 & 36873.66 & 63966.54 \\
    \bottomrule
    \end{tabular*}
  \end{threeparttable}
\end{table}

\subsection{西安参数下的隔离率阈值控制}
西安市 2021 年末常住人口约 1316.30 万~\cite{XianStat2021}；《2022 年陕西省卫生健康统计年鉴》给出 2021 年西安市卫生机构床位 79426 张~\cite{ShaanxiHealth2022}，约为 $0.006N$。考虑疫情期间不可能将全部床位用于收治，按约三分之一可调配估算，取代表性阈值
\begin{equation}
\eta=0.002N=26{,}326.
\end{equation}
按第~\ref{sec:s1}节，常规控制下临界易感人数 $S_c=\gamma N/[\beta c_0(1-q_0)]=2{,}974{,}125.41$；数值求得达到阈值时刻 $t_1=16.90$、启动点 $S^*=13{,}020{,}440.75$。阈值控制期由理论开环控制 $q_c(t)=1-\gamma N/[\beta c_0 S(t)]$（其中 $S(t)=\bar S+(S^*-\bar S)e^{-c_0\eta(t-t_1)/N}$，$\bar S=\gamma N(1-\beta)/(\beta c_0)$）维持 $I(t)\equiv\eta$，解除时刻
\[
t_2=t_1+\frac{N}{c_0\eta}\ln\frac{S^*-\bar S}{S_c-\bar S}\approx101.97,
\]
控制时长约 $85.07$ 天；启动隔离率 $q_c(t_1)=0.8454$，终点 $q_c(t_2)=q_0=0.3230$。阈值控制阶段最大数值偏差约为 $2.89\times10^{-8}$。

\subsection{比较指标与三策略数值结果}
三类策略均计算至各自的动态清零时刻，即 $I(t)$ 超过 $1$ 后再次下降至 $1$ 的时刻。成本采用第~\ref{sec:cost}~节的定义，权重为 $w_c=1,w_q=2$：TDINN 从初始时刻积分至清零，阈值控制仅在 $[t_1,t_2]$ 内产生额外成本，常规控制的额外成本为零。

图~\ref{fig:xian_panels}--\ref{fig:xian_Re} 分别给出三类策略下的感染与控制轨迹、累计感染和有效再生数，主要指标列于表~\ref{tab:xian_summary}。

\begin{figure}[htbp]
    \centering
    \includegraphics[width=0.7\textwidth]{xian_control_comparison_panels.pdf}
    \caption{$\eta=0.002N$ 时三类策略的比较。(a) 非隔离感染者人数；(b)、(c) 每日新增社区和隔离感染；(d)、(e) 接触率和隔离率。(b)、(c) 内嵌图中的点为观测数据，竖线表示观测结束及阈值控制的启动、结束时刻。}
    \label{fig:xian_panels}
\end{figure}

\begin{figure}[htbp]
    \centering
    \includegraphics[width=0.7\textwidth]{xian_control_cumulative.pdf}
    \caption{三类策略下的累计感染。(a) 社区累计感染 $I_{\rm cum}$；(b) 隔离累计感染 $I_{q_{\rm cum}}$；(c) 总累计感染 $\Itcum$。}
    \label{fig:xian_cumulative}
\end{figure}

\begin{figure}[htbp]
    \centering
    \includegraphics[width=0.7\textwidth]{xian_effective_reproduction_number.pdf}
    \caption{三类策略下有效再生数 $R_e(t)$ 的变化。水平虚线为 $R_e=1$，阈值控制在阈值控制期维持这一临界水平。}
    \label{fig:xian_Re}
\end{figure}

\begin{table}[htbp]
  \centering
  \small
  \sisetup{group-digits=false}
  \setlength{\tabcolsep}{3.5pt}
  \renewcommand{\arraystretch}{1.25}
  \begin{threeparttable}
    \caption{$\eta=0.002N$ 下三类控制策略的数值比较}
    \label{tab:xian_summary}
    \begin{tabular*}{\textwidth}{@{\extracolsep{\fill}}l*{3}{S[table-format=7.2]}@{}}
    \toprule
    比较量 & {TDINN控制} & {情景一阈值控制} & {常规控制} \\
    \midrule
    感染峰值 $\max I$ & 152 & 26326 & 1377550 \\
    启动时间 $t_1$（d） & 0.00 & 16.90 & 0.00 \\
    解除时间 $t_2$（d） & 45.27 & 101.97 & 0.00 \\
    控制时长 $\Delta t$（d） & 45.27 & 85.07 & 0.00 \\
    清零时间 $t_{\rm end}$（d） & 45.27 & 258.11 & 75.58 \\
    总累计感染 $\Itcum$ & 2097 & 2377943 & 4587183 \\
    \addlinespace[3pt]
    接触率分项成本 $J_c$ & 19.15 & 0.00 & 0.00 \\
    隔离率分项成本 $J_q$ & 25.15 & 36.80 & 0.00 \\
    二次加权总成本 $J$ & 49.35 & 40.77 & 0.00 \\
    \bottomrule
    \end{tabular*}
    \begin{tablenotes}[flushleft]
      \footnotesize
      \item 注：表内 $\max I$、$\Itcum$ 为直接模拟取整值。
    \end{tablenotes}
  \end{threeparttable}
\end{table}

采用全市人口及 $\eta=0.002N$ 时，阈值控制将社区感染峰值限制为 $26{,}326$，阈值控制持续约 $85.07$ 天，动态清零时间约为 $258.11$ 天，总累计感染约为 $2.38\times10^6$。TDINN 重构的峰值和累计感染分别约为 $152$ 和 $2097$。阈值控制的二次加权成本为 $40.77$，低于 TDINN 的 $49.35$，但感染峰值、累计感染和清零时间均较大。后文采用更精确的 TDINN 参照值 $I_{\rm peak}^{\rm T}=151.90$、$\Itcum^{\rm T}=2096.76$ 和 $J^{\rm T}=49.35$。

有效再生数为 $R_e(t)=\beta c(t)(1-q(t))S(t)/(\gamma N)$。常规控制下 $R_e(0)\approx4.43$，感染人数初期快速增长；TDINN 同时降低接触率并提高隔离率，使 $R_e$ 较快降至 $1$ 以下；阈值控制则在阈值控制期保持 $R_e=1$，待 $S$ 下降至 $S_c$ 后恢复常规隔离率，感染人数才开始下降。因此，较低的加权成本在这里伴随着更长的感染持续时间和更高的累计感染。

\subsection{阈值 \texorpdfstring{$\eta$}{eta} 的敏感性}
取 $\eta/N\in\{0.0001,0.0002,0.0005,0.001,0.0015,0.002,0.003,0.004,0.005,0.0075,0.01\}$，结果见表~\ref{tab:xian_eta_sens} 和图~\ref{fig:xian_eta_sens}。在这一范围内，有
\[
\eta\downarrow\ \Rightarrow\ t_1\downarrow,\qquad \Delta t\uparrow,\qquad J\uparrow,\qquad t_{\rm end}\uparrow.
\]
当 $\eta/N$ 从 $0.01$ 降至 $0.0001$ 时，启动时间仅由 $18.55$ 天提前至 $13.93$ 天，控制持续时间却由 $16.61$ 天增加至 $1710.66$ 天，动态清零时间达到 $2209.55$ 天。降低阈值减慢了阈值控制期的易感者减少过程，其对持续时间的影响远大于对启动时间的影响。

\begin{table}[htbp]
  \centering
  \small
  \sisetup{group-digits=false}
  \setlength{\tabcolsep}{3.5pt}
  \renewcommand{\arraystretch}{1.25}
  \begin{threeparttable}
    \caption{不同医疗阈值 $\eta$ 下隔离率阈值控制的数值指标}
    \label{tab:xian_eta_sens}
    \begin{tabular*}{\textwidth}{@{\extracolsep{\fill}}S[table-format=6.0]S[table-format=1.4]S[table-format=2.2]S[table-format=4.2]S[table-format=4.2]S[table-format=4.2]S[table-format=1.4]S[table-format=3.2]S[table-format=7.0]@{}}
    \toprule
    \multicolumn{2}{c}{阈值设置} & \multicolumn{4}{c}{时间量（d）} & \multicolumn{2}{c}{隔离强度与成本} & {总累计感染} \\
    \cmidrule(lr){1-2}\cmidrule(lr){3-6}\cmidrule(lr){7-8}
    {$\eta$} & {$\eta/N$} & {$t_1$} & {$t_2$} & {$\Delta t$} & {$t_{\rm end}$} & {$q_c(t_1)$} & {$J$} & {$\Itcum$} \\
    \midrule
    1316 & 0.0001 & 13.93 & 1724.59 & 1710.66 & 2209.55 & 0.8470 & 826.41 & 2155649 \\
    2633 & 0.0002 & 14.61 & 869.70 & 855.09 & 1243.16 & 0.8469 & 412.91 & 2182123 \\
    6582 & 0.0005 & 15.52 & 357.27 & 341.75 & 620.77 & 0.8466 & 164.82 & 2234666 \\
    13163 & 0.001 & 16.21 & 186.84 & 170.63 & 389.32 & 0.8462 & 82.12 & 2293939 \\
    19744 & 0.0015 & 16.61 & 130.20 & 113.59 & 304.01 & 0.8458 & 54.55 & 2339486 \\
    26326 & 0.002 & 16.90 & 101.97 & 85.07 & 258.11 & 0.8454 & 40.77 & 2377943 \\
    39489 & 0.003 & 17.31 & 73.86 & 56.55 & 208.39 & 0.8445 & 26.98 & 2442608 \\
    52652 & 0.004 & 17.60 & 59.89 & 42.29 & 181.18 & 0.8436 & 20.09 & 2497306 \\
    65815 & 0.005 & 17.83 & 51.56 & 33.73 & 163.66 & 0.8428 & 15.95 & 2545663 \\
    98722 & 0.0075 & 18.25 & 40.57 & 22.32 & 138.11 & 0.8405 & 10.44 & 2649240 \\
    131630 & 0.01 & 18.55 & 35.16 & 16.61 & 123.89 & 0.8382 & 7.68 & 2737382 \\
    \bottomrule
    \end{tabular*}
  \end{threeparttable}
\end{table}

\begin{figure}[htbp]
    \centering
    \includegraphics[width=0.7\textwidth]{xian_eta_sensitivity.pdf}
    \caption{西安参数下控制指标随 $\eta/N$ 的变化。(a) 启动隔离率；(b) 加权成本；(c) 累计感染；(d) 控制持续时间与动态清零时间。水平线为 TDINN 的对应参照值。}
    \label{fig:xian_eta_sens}
\end{figure}

\FloatBarrier
\subsection{接触率 \texorpdfstring{$c_0$}{c0} 与阈值 \texorpdfstring{$\eta$}{eta} 的联合变化}
图~\ref{fig:xian_heatmaps} 给出 $(\eta,c_0)$ 平面上的控制持续时间、清零时间、累计感染和成本。两条参照线分别为 $\eta=0.002N=26326$ 和 $\eta=I_{\rm peak}^{\rm T}=151.90$。

在西安全市人口设定（$N=13\,163\,000$，$c_0=12.8872$）下，由控制时长闭式
\[
\Delta t=\frac{N}{c_0\eta}\ln\frac{S^{*}-\bar S}{S_c-\bar S}
\]
有$\eta=0.002N=26326$ 时 $\Delta t\approx85$ 天；$\eta=151.90$ 时，$S^{*}\approx1.3162\times10^{7}$、对数因子约为 $2.205$，从而
\[
\Delta t(\eta=151.90)\approx1.48\times10^{4}\ \text{d}\approx 40.6\ \text{年}.
\]
这一时长是模型在固定参数和全市人口假设下的外推结果。由时间尺度 $N/(c_0\eta)$ 可见，较大的易感人群和较低的感染阈值会延长阈值控制期。在观测期累计病例约为 $2050$ 的拟合中，$S(t)\approx N$；但阈值策略需要持续至 $S=S_c$，因此其时间估计对人口规模尤为敏感。下一节考察不同有效混合人口下的结果。

\begin{figure}[htbp]
    \centering
    \includegraphics[width=0.8\textwidth]{xian_heatmaps.pdf}
    \caption{西安参数下 $(\eta,c_0)$ 平面上的控制指标。(a) 控制持续时间；(b) 动态清零时间；(c) 累计感染；(d) 加权成本。竖直参照线为 $\eta=151.90$ 和 $\eta=26326$，等值线表示各指标的水平。}
    \label{fig:xian_heatmaps}
\end{figure}


\section{规模不变性与不同有效人口下的策略比较}
\label{sec:dominance}

西安此次疫情的有效混合人口 $N_{\rm eff}$ 未知。本节沿用第~\ref{sec:xian}~节的参数，改变 $N_{\rm eff}$，考察阈值控制满足 TDINN 峰值、成本及时间参照值的参数范围。TDINN 参照固定在全市人口下的拟合结果，阈值控制和常规控制则在各个人口假设下分别计算。

\subsection{设定与比较基准}
\label{sec:dom:setup}

记
\begin{equation}
  \theta:=\frac{\eta}{N}\quad(\text{阈值比例}),\qquad
  i_{\max}^{no}:=\frac{I_{\max}^{no}}{N}\quad(\text{常规控制峰值比例}),
  \label{eq:dom:ratios}
\end{equation}
其中 $I_{\max}^{no}$ 由式~\eqref{eq:s1:Imax-no} 给出。在归一化初值 $s_0=S_0/N$ 与 $i_0=I_0/N$ 固定时，
$S_c,\bar S,I_{\max}^{no}$ 均随 $N$ 同比例缩放，故 $i_{\max}^{no}$ 不依赖 $N$，
只依赖 $\beta,\gamma,c_0,q_0,s_0,i_0$。本节记归一化状态 $s=S/N$、$i=I/N$。

数值计算固定绝对初值 $I_0=1.00663\times10^{-3}$，不随 $N_{\rm eff}$ 重新拟合，故 $i_0=I_0/N$ 随人口规模改变。以下理论中的规模不变性则以固定归一化初值为前提，两者需加以区分。

由首次积分，$i_0$ 通过 $\theta-i_0$ 和 $s_0=1-i_0$ 影响启动点。这一影响的数量级估计为 $O(i_0/\theta)$；在 $N\ge1.06\times10^4$、$\theta\ge1.14\times10^{-3}$ 时，$i_0/\theta\le8\times10^{-5}$。数值计算所得 $\Delta t$ 和 $J$ 的跨人口相对极差分别为 $2.4\times10^{-7}$ 和 $5.9\times10^{-7}$；当 $i_0$ 在 $10^{-11}$ 至 $10^{-7}$ 内变化时，$\theta_{\rm cost}$ 和 $N^\ast$ 在七位有效数字内未显示变化。

TDINN 的四个参照量见表~\ref{tab:dom:tdinn}。社区峰值、累计感染和清零时间由拟合轨迹重构，成本则由拟合控制函数按 $w_c=1,w_q=2$ 积分得到；这些量均不随本节的 $N_{\rm eff}$ 改变。

\begin{table}[htbp]
  \centering
  \small
  \sisetup{group-digits=false}
  \setlength{\tabcolsep}{3.5pt}
  \renewcommand{\arraystretch}{1.25}
  \begin{threeparttable}
    \caption{TDINN 比较基准：西安全市尺度下的四个参照量}
    \label{tab:dom:tdinn}
\begin{tabularx}{\textwidth}{@{}>{\raggedright\arraybackslash}X>{\centering\arraybackslash}X S[table-format=4.2]@{}}
    \toprule
    参照量 & 符号 & {数值} \\
    \midrule
    社区峰值 & $I_{\rm peak}^{\rm T}$ & 151.90 \\
    总累计感染 & $\Itcum^{\rm T}$ & 2096.76 \\
    清零时刻（d） & $t_{\rm end}^{\rm T}$ & 45.27 \\
    二次加权成本 & $J^{\rm T}$ & 49.35 \\
    \bottomrule
  \end{tabularx}
    \begin{tablenotes}[flushleft]
      \footnotesize
      \item 注：前三项由数据重构得到；$J^{\rm T}$ 由式~\eqref{eq:cost:J} 对拟合控制函数积分得到。
    \end{tablenotes}
  \end{threeparttable}
\end{table}

\subsection{规模不变性与单调性}
\label{sec:dom:theory}

\begin{lemma}\label{lem:dom:scaling}
固定归一化初值 $S_0/N,I_0/N$ 及传播和常规控制参数，在满足阈值触发条件时，控制持续时间 $\Delta t$、二次加权成本 $J$、
最大隔离率 $q_{\max}$ 以及峰值比例 $\eta/N$ 在 $(N,\eta)$ 中都只与比例 $\theta=\eta/N$ 有关；
特别地 $s^*:=S^\ast/N$ 只与 $\theta$ 有关。
\end{lemma}
\begin{proof}
记 $s_c:=S_c/N=\gamma/[\beta c_0(1-q_0)]$、$\bar s:=\bar S/N=\gamma(1-\beta)/(\beta c_0)$，
二者不含 $N$。将式~\eqref{eq:s1:Sstar-lambert} 的 Lambert $W$ 函数的自变量改写：由
$\rho_1=\gamma N/\{c_0[\beta+q_0(1-\beta)]\}$、
$\frac{\beta_2^0/\beta_1^0}{\rho_1}=1/S_c$ 与
$C_0=I_0+\tfrac{\beta_2^0}{\beta_1^0}S_0-\rho_1\ln S_0$，
\[
  -\frac{C_0-\eta}{\rho_1}
  = \ln S_0-\frac{S_0}{S_c}-\frac{I_0-\eta}{\rho_1}
  = \ln N+\left[
  \ln\frac{S_0}{N}-\frac{S_0/N}{s_c}
  -\frac{c_0[\beta+q_0(1-\beta)]}{\gamma}
  \left(\frac{I_0}{N}-\theta\right)
  \right].
\]
于是该自变量可写为
\[
-\frac{1}{s_c}
\exp\left[
\ln\frac{S_0}{N}-\frac{S_0/N}{s_c}
-\frac{c_0[\beta+q_0(1-\beta)]}{\gamma}
\left(\frac{I_0}{N}-\theta\right)
\right],
\]
其中 $e^{\ln N}$ 已与 $S_c=Ns_c$ 中的 $N$ 抵消。该式不含 $N$，
仅依赖 $\theta$ 与归一化初值，故
$s^*=-s_c W_{-1}(\cdot)$ 也只与 $\theta$ 有关。将其代入
$\Delta t=\frac{N}{c_0\eta}\ln\frac{S^\ast-\bar S}{S_c-\bar S}
 =\frac{1}{c_0\theta}\ln\frac{s^*-\bar s}{s_c-\bar s}$、
$q_{\max}=1-\frac{\gamma N}{\beta c_0 S^\ast}=1-\frac{\gamma}{\beta c_0\,s^*}$，
以及（由 $\dot S=-\frac{c_0\eta}{N}(S-\bar S)$、$dt=-\frac{N}{c_0\eta}\frac{dS}{S-\bar S}$）
\begin{equation}
  J=\frac{2N}{c_0\eta(1-q_0)^2}\int_{S_c}^{S^\ast}\frac{\big(q_c(S)-q_0\big)^2}{S-\bar S}\,dS
   =\frac{2}{c_0\theta(1-q_0)^2}\int_{s_c}^{s^*}\frac{\big(q_c(N s)-q_0\big)^2}{s-\bar s}\,ds,
  \label{eq:dom:Jtheta}
\end{equation}
其中 $q_c(N s)=1-\gamma/(\beta c_0 s)$ 不含 $N$。在其他参数和归一化初值固定时，右端均只通过 $\theta$ 依赖于 $N,\eta$。
\end{proof}

固定绝对初值的数值计算不完全满足上述引理的前提，其近似程度由第~\ref{sec:dom:setup}~小节的数值结果说明。

\begin{lemma}\label{lem:dom:mono}
固定归一化初值和其余参数，$\Delta t$ 与 $J$ 在 $\theta\in(i_0,i_{\max}^{no})$ 上均严格单调递减。
此外，沿任意满足
\[
0<i_0<\theta,\qquad
i_0\to0,\qquad
\theta\to0
\]
的可行极限序列，若 $s_0>s_c$，
则
\[
\Delta t\to+\infty,
\qquad
J\to+\infty.
\]
\end{lemma}
\begin{proof}
由引理~\ref{lem:dom:scaling}，$\Delta t,J$ 是 $\theta$ 的函数；固定 $N$ 时 $\theta$ 与 $\eta$ 同增，
故其对 $\theta$ 的单调性等同于对 $\eta$ 的单调性。推论~\ref{cor:s1:duration-sensitivity} 已给出
$\frac{\partial \Delta t}{\partial \eta}<0$，引理~\ref{lem:s1:Sstar-sensitivity} 给出 $\frac{\partial S^\ast}{\partial \eta}<0$。对 $J$，式~\eqref{eq:dom:Jtheta} 中的系数 $1/(c_0\theta)$ 随 $\theta$ 严格递减，积分上限 $s^*$ 也随 $\theta$ 减小。由于被积函数非负且在积分区间内部为正，故 $J$ 严格递减。

对于低阈值极限，
沿任意满足
$0<i_0<\theta$、
$i_0\to0$、
$\theta\to0$
的可行序列，阈值方程的连续性给出
\[
  s^\ast\longrightarrow s_0.
\]
由于假设 $s_0>s_c>\bar s$，
有
\[
  \ln\frac{s^\ast-\bar s}{s_c-\bar s}
  \longrightarrow
  \ln\frac{s_0-\bar s}{s_c-\bar s}
  >0.
\]
因此
\[
  \Delta t
  =
  \frac{1}{c_0\theta}
  \ln\frac{s^\ast-\bar s}{s_c-\bar s}
  \longrightarrow+\infty.
\]

对成本，式~\eqref{eq:dom:Jtheta} 中的积分满足
\[
\int_{s_c}^{s^\ast}
\frac{\bigl(q_c(Ns)-q_0\bigr)^2}
     {s-\bar s}\,ds
\longrightarrow
K_0
:=
\int_{s_c}^{s_0}
\frac{\bigl(q_c(Ns)-q_0\bigr)^2}
     {s-\bar s}\,ds.
\]
由
$q_c(Ns_c)=q_0$，
且 $q_c(Ns)>q_0$ 对所有
$s\in(s_c,s_0]$ 成立，
同时 $s-\bar s>0$，
故 $K_0>0$。于是
\[
  J(\theta)
  \sim
  \frac{2K_0}
       {c_0(1-q_0)^2}\,
  \frac{1}{\theta}
  \longrightarrow+\infty.
\]
\end{proof}

由引理~\ref{lem:dom:mono}，若 $J^{\rm T}$、$T_{\max}$ 落在相应值域内，则下列阈值唯一存在：
\begin{equation}
  J\big(\theta_{\rm cost}\big)=J^{\rm T},\qquad
  \Delta t\big(\theta_{\rm dur}(T_{\max})\big)=T_{\max}.
  \label{eq:dom:thetadefs}
\end{equation}
为简化记号，以下记
\begin{equation}
  \theta_{\rm bind}(T_{\max})
  :=
  \max\!\bigl(
    \theta_{\rm cost},
    \theta_{\rm dur}(T_{\max})
  \bigr),
  \label{eq:dom:theta-bind}
\end{equation}
在不引起混淆时简写为 $\theta_{\rm bind}$。

\subsection{有效人口的下界与占优区域}
\label{sec:dom:region}

在无回流模型下，可由累计感染量得到不依赖于控制函数具体形式的有效人口下界。

\begin{proposition}\label{prop:dom:floor}
对于本节的西安重构轨迹，记 $I_{\rm cum}^{\rm T},I_{q_{\rm cum}}^{\rm T}$ 为累计社区感染和累计隔离感染，则
\begin{equation}
  N_{\rm eff}\ \ge\ N_{\rm floor}
  :=I_{\rm cum}^{\rm T}+\frac{I_{q_{\rm cum}}^{\rm T}}{\beta}
  =1.06\times10^{4}.
  \label{eq:dom:Nfloor}
\end{equation}
\end{proposition}
\begin{proof}
由式~\eqref{eq:model:full}，
\[
  \dot S_q=\frac{(1-\beta)c(t)q(t)}{N}SI,
  \qquad
  \dot I_{q_{\rm cum}}=\frac{\beta c(t)q(t)}{N}SI
  \quad\Longrightarrow\quad
  \dot S_q=\frac{1-\beta}{\beta}\,\dot I_{q_{\rm cum}},
\]
而 $S_q$ 只有流入、无回流。易感者的三条去向为社区感染、隔离感染与被追踪隔离的未感染者，
故该次疫情消耗的易感者总数为
\begin{equation}
  I_{\rm cum}^{\rm T}+I_{q_{\rm cum}}^{\rm T}+S_q^{\rm T}
  =I_{\rm cum}^{\rm T}+\frac{I_{q_{\rm cum}}^{\rm T}}{\beta}
  =605.40+\frac{1491.36}{0.1498}
  =1.0561\times10^{4}.
  \label{eq:dom:Sconsumed}
\end{equation}
上述关系不依赖于 $c(t),q(t)$ 的具体时间形式。
再由 $S_0\le N$ 即得式~\eqref{eq:dom:Nfloor}。
\end{proof}

式~\eqref{eq:dom:Nfloor} 是本节区间的下端。%
\footnote{$N_{\rm floor}$ 依赖 $\beta$：$\beta=0.20$ 时为 $8062$，$0.25$ 时为 $6571$，
$0.30$ 时为 $5577$。附录~\ref{app:dom:numerics} 说明区间的宽度对 $\beta$ 稳健。}
仅计累计感染时，相应人数为 $\Itcum^{\rm T}=2096.76$；计入隔离未感染者后，$N_{\rm floor}$ 约为该值的 $5.04$ 倍。由于总累计感染不超过人口总数，$N<\Itcum^{\rm T}$ 时累计比较条件自动满足，但这一人口范围低于模型给出的 $N_{\rm floor}$。接近 $N_{\rm floor}$ 意味着社区易感者大部分离开 $S$ 仓室，见第~\ref{sec:dom:limits}~小节。

下面由峰值、成本和时长条件确定有效人口上界。

\begin{proposition}\label{prop:dom:region}
在引理~\ref{lem:dom:mono} 的条件下，设 $\theta_{\rm cost}$ 和 $\theta_{\rm dur}(T_{\max})$ 按式~\eqref{eq:dom:thetadefs} 存在。隔离率阈值控制同时满足
(i) $I_{\rm peak}^{\rm thr}=\eta\le I_{\rm peak}^{\rm T}$、
(ii) $J\le J^{\rm T}$、
(iii) $\Delta t\le T_{\max}$，
当且仅当
\begin{equation}
  \theta_{\rm bind}(T_{\max})
  \ \le\
  \frac{\eta}{N}
  \ \le\
  \frac{I_{\rm peak}^{\rm T}}{N},
  \qquad\text{且}\qquad
  \frac{\eta}{N}<i_{\max}^{no}.
  \label{eq:dom:band}
\end{equation}

式~\eqref{eq:dom:band} 所给区间对 $\eta$ 非空，
当且仅当
\begin{equation}
  \theta_{\rm bind}(T_{\max})<i_{\max}^{no}
  \qquad\text{且}\qquad
  N\le N^\ast(T_{\max})
  :=
  \frac{I_{\rm peak}^{\rm T}}
       {\theta_{\rm bind}(T_{\max})}.
  \label{eq:dom:Nstar}
\end{equation}
特别地，若
$\theta_{\rm bind}(T_{\max})\ge i_{\max}^{no}$，
则该占优区间对所有 $N$ 均为空。
\end{proposition}
\begin{proof}
条件 (i) 等价于
$\eta/N\le I_{\rm peak}^{\rm T}/N$。
由引理~\ref{lem:dom:mono} 与式~\eqref{eq:dom:thetadefs}，
条件 (ii) 等价于 $\eta/N\ge\theta_{\rm cost}$，
条件 (iii) 等价于 $\eta/N\ge\theta_{\rm dur}(T_{\max})$。
再与控制持续时间为正的触发条件 $\eta/N<i_{\max}^{no}$ 取交，即得式~\eqref{eq:dom:band}。
触发条件取严格不等式以排除 $\Delta t=0$；三个比较条件允许等号。

该区间非空当且仅当
$\theta_{\rm bind}(T_{\max})\le I_{\rm peak}^{\rm T}/N$
与
$\theta_{\rm bind}(T_{\max})<i_{\max}^{no}$
同时成立。由于 $\theta_{\rm bind}(T_{\max})>0$，前一个不等式等价于
$N\le I_{\rm peak}^{\rm T}/\theta_{\rm bind}(T_{\max})=N^\ast(T_{\max})$，
即得式~\eqref{eq:dom:Nstar}；
若 $\theta_{\rm bind}(T_{\max})\ge i_{\max}^{no}$，则区间对任意 $N$ 均为空。
\end{proof}

\begin{remark}
这里的占优仅指满足所列指标的比较条件；其中峰值条件 $\eta\le I_{\rm peak}^{\rm T}$ 是对阈值的预设要求，不代表该策略在其他指标上也较低。
\end{remark}

\begin{corollaryn}\label{cor:ceiling}
由引理~\ref{lem:dom:mono}，$\theta_{\rm dur}(T_{\max})$ 随 $T_{\max}$ 增大而单调递减，
且 $\theta_{\rm dur}(T_{\max})\to0\ (T_{\max}\to\infty)$，因此
\[
  \theta_{\rm bind}(T_{\max})
  =
  \max\!\bigl(
    \theta_{\rm cost},
    \theta_{\rm dur}(T_{\max})
  \bigr)
  \longrightarrow
  \theta_{\rm cost}.
\]
若 $\theta_{\rm cost}<i_{\max}^{no}$，则在满足 $\theta_{\rm bind}(T_{\max})<i_{\max}^{no}$
的参数范围内，$N^\ast(T_{\max})=I_{\rm peak}^{\rm T}/\theta_{\rm bind}(T_{\max})$
随 $T_{\max}$ 单调不减，并且
\[
  N^\ast(T_{\max})
  \longrightarrow
  N^\ast_\infty
  :=
  \frac{I_{\rm peak}^{\rm T}}
       {\theta_{\rm cost}}.
\]
故即使允许任意长的阈值控制期，有效人口上界仍被成本约束限制在 $N^\ast_\infty$。
若 $\theta_{\rm cost}\ge i_{\max}^{no}$，则即使令 $T_{\max}\to\infty$，
命题~\ref{prop:dom:region} 的占优区仍对所有 $N$ 为空。
\end{corollaryn}

\subsection{总累计与清零时刻的附加约束}
\label{sec:dom:extensive}

以下在峰值、成本和时长条件之外，加入累计感染与动态清零时间的比较。

\begin{corollaryn}\label{cor:cum}
在命题~\ref{prop:dom:region} 的占优区域上再加约束
(iv) $\Itcum^{\rm thr}\le\Itcum^{\rm T}$。由定理~\ref{thm:s1:Itcum}，
到清零的总累计感染为
\[
  \Itcum^{\rm thr}=N\,h(\theta,N),\qquad \theta=\eta/N,
\]
其退出段因绝对清零判据 $I=1$ 带有弱 $1/N$ 依赖，故 $h=h(\theta,N)$。
累计感染的比较边界由下式给出
\begin{equation}
  N\,h(\theta,N)=\Itcum^{\rm T}
  \label{eq:dom:cum-feasible}
\end{equation}
在 $(N,\eta)$ 平面上表现为曲线，而不是过原点的直线，占优侧（$\Itcum^{\rm thr}\le\Itcum^{\rm T}$）
为其小 $N$ 一侧。因此累计约束把占优区域整体收缩。

若对每个固定 $N$，$h(\theta,N)$ 关于 $\theta$ 单调不减，则累计感染约束从较大 $\theta$ 的一侧缩小比较区域，
而占优区域的下边缘 $\theta=\theta_{\rm bind}(T_{\max})$ 最小，因此，累计感染边界与该下界的交点决定相应的有效人口上限。
记 $F_{T_{\max}}(N):=N\,h\!\bigl(\theta_{\rm bind}(T_{\max}),N\bigr)$，该交点的有效人口
$N_{\rm cum}^\ast(T_{\max})$ 由
\begin{equation}
  F_{T_{\max}}\!\bigl(N_{\rm cum}^\ast(T_{\max})\bigr)=\Itcum^{\rm T}
  \label{eq:dom:Ncumstar}
\end{equation}
确定。由上述单调性与交点条件，纳入累计约束后的占优区域对 $\eta$ 非空，当且仅当
\[
N\le\min\{N^\ast(T_{\max}),\,N_{\rm cum}^\ast(T_{\max})\}.
\]
这一条件由累计约束与式~\eqref{eq:dom:Nstar} 取交得到。
两个上界的大小由具体参数决定。西安参数下的计算结果见第~\ref{sec:dom:xian}~小节。
\end{corollaryn}

\begin{corollaryn}\label{cor:clr}
阈值控制到动态清零的时刻 $t_{\rm end}$ 由定理~\ref{thm:s1:tend} 给出。由于清零条件为 $I=1$，即 $i=1/N$，退出段还依赖于绝对人口规模，故 $t_{\rm end}=t_{\rm end}(\theta,N)$。本节数值结果中，固定 $\theta$ 时清零时间随 $N$ 增大。若再要求
(v) $t_{\rm end}^{\rm thr}\le t_{\rm end}^{\rm T}$，
则在占优区域内该约束给出一条 $(N,\eta)$ 平面上的曲线边界
\begin{equation}
  t_{\rm end}\big(\eta/N,\,N\big)=t_{\rm end}^{\rm T},
  \label{eq:dom:clrarc}
\end{equation}
其占优侧为小 $N$ 一侧；记其与给定 $\eta$ 的交点为 $N_{\rm clr}(\eta)$。由式~\eqref{eq:dom:Nfloor}，
与人口下界共同确定的区间为 $N\in[\,N_{\rm floor},\,N_{\rm clr}(\eta)\,]$。本文参数下
$\max_\eta N_{\rm clr}(\eta)=5166<N_{\rm floor}=1.06\times10^{4}$，
故该区间为空集：不存在同时满足下界与清零不劣的有效人口（第~\ref{sec:dom:xian}~小节）。
\end{corollaryn}

\begin{remarkn}\label{rem:dom:dichotomy}
在固定归一化初值和其余参数时，$\Delta t,J,q_{\max}$ 仅依赖于 $\theta=\eta/N$，其等值条件可写成 $\eta=\theta N$。绝对峰值约束 $\eta\le I_{\rm peak}^{\rm T}$ 以及累计感染、清零时间的条件仍涉及人口规模。因此，相同 $\theta$ 并不意味着所有比较条件相同。
\end{remarkn}

\subsection{规模不变性的数值验证}
\label{sec:scaling}

采用西安参数和固定绝对初值，分别固定阈值比例及绝对阈值，比较不同 $N_{\rm eff}$ 下的控制结果。

固定 $\theta=0.002$，取 $N_{\rm eff}\in\{5\times10^4,\dots,1.316\times10^7\}$。六个人口规模下的控制持续时间均约为 $85.07$ 天，相对极差为 $2.4\times10^{-7}$。以 $t-t_1$ 为时间变量时，归一化感染轨迹和隔离率近似重合，见图~\ref{fig:neff_collapse}。启动时间则由 $11.39$ 天增加至 $16.90$ 天，这是固定绝对初值、初始感染比例随 $N$ 变化的结果。

固定绝对阈值 $\eta=100$，把 $N_{\rm eff}$ 从 $1.316\times10^7$ 降到 $5\times10^{4}$，
则 $\theta=\eta/N_{\rm eff}$ 增大，控制持续时间由 $\Delta t\approx2.25\times10^4$ 天缩短到约
$85$ 天（图~\ref{fig:neff_time}）。这是降低 $N_{\rm eff}$ 通过增大 $\theta$ 缩短阈值控制期的数值体现，
与提高 $\eta$ 同向。

两组结果中，控制持续时间主要随 $\theta$ 变化：
\[
\Delta t\approx\frac{1}{c_0\theta}\ln\frac{s^*-\bar s}{s_c-\bar s}\approx\frac{\text{const}}{c_0\theta},
\]
其中对数因子在 $\theta\in[5\times10^{-4},8\times10^{-3}]$ 内仅由 $2.20$ 变到 $2.15$，
故 $\Delta t$ 近似与 $\theta$ 成反比。

\begin{figure}[htbp]
    \centering
    \includegraphics[width=0.7\textwidth]{neff_scaling_collapse.pdf}
    \caption{固定绝对初值和 $\theta=0.002$ 时，不同有效人口下的归一化感染轨迹与隔离率，时间变量为 $t-t_1$。空心圆表示各条轨迹达到 $I=1$ 的时刻。}
    \label{fig:neff_collapse}
\end{figure}

\begin{figure}[htbp]
    \centering
    \includegraphics[width=0.95\textwidth]{neff_time_metrics.pdf}
    \caption{固定 $\eta=100$ 时，启动时间、控制持续时间和动态清零时间随有效人口的变化。图中 $T_{\rm clear}$ 表示动态清零时间。}
    \label{fig:neff_time}
\end{figure}

固定 $\beta=0.1498$，在所考察的 $\eta\in\{80,100,150\}$、$N_{\rm eff}\in\{4,5,6\}\times10^4$ 参数组中，拐点处的隔离率均约为 $0.4119$，与理论值的最大绝对误差约为 $8\times10^{-9}$。这与 $\qinf$ 仅依赖于 $\beta$ 的表达式一致。当 $q_0=0.3230$ 时，条件 $\qinf=q_0$ 给出 $\beta_{\max}=0.261448$；超过该值后，阈值控制区间内不再有内部拐点。

\subsection{西安参数下的占优区域}
\label{sec:dom:xian}

取 $\beta=0.1498,\gamma=0.2953,c_0=12.8872,q_0=0.3230$，$S_0/N\approx1$，参照量见表~\ref{tab:dom:tdinn}。由式~\eqref{eq:s1:Imax-no}，$i_{\max}^{no}\approx0.1047$。表~\ref{tab:dom:thresholds} 给出阈值比例和临界有效人口；在所考察的 $T_{\max}=45,60,90,150$ 天及不限制时长的设定下，均有 $\theta_{\rm bind}<i_{\max}^{no}$。
当 $T_{\max}\approx103$ 天时，$\theta_{\rm dur}=\theta_{\rm cost}$；更短的允许时长使时长条件起约束作用，更长的允许时长则由成本条件决定阈值下限。

\begin{table}[htbp]
  \centering
  \small
  \sisetup{group-digits=false}
  \setlength{\tabcolsep}{3.5pt}
  \renewcommand{\arraystretch}{1.25}
  \begin{threeparttable}
    \caption{西安参数下的阈值比例与临界有效人口}
    \label{tab:dom:thresholds}
\begin{tabularx}{\textwidth}{@{}>{\raggedright\arraybackslash}p{0.18\textwidth}>{\centering\arraybackslash}X>{\centering\arraybackslash}p{0.23\textwidth}>{\raggedright\arraybackslash}p{0.21\textwidth}@{}}
    \toprule
    类别 & 变量 & 数值 & 依据 \\
    \midrule
    \multirow{4}{*}{阈值比例} & $\theta_{\rm cost}$ & $1.656\times10^{-3}$ & 式~\eqref{eq:dom:thetadefs} \\
    & $\theta_{\rm dur}(45\,\mathrm d)$ & $3.762\times10^{-3}$ & 式~\eqref{eq:dom:thetadefs} \\
    & $\theta_{\rm dur}(90\,\mathrm d)$ & $1.891\times10^{-3}$ & 式~\eqref{eq:dom:thetadefs} \\
    & $\theta_{\rm dur}(150\,\mathrm d)$ & $1.137\times10^{-3}$ & 式~\eqref{eq:dom:thetadefs} \\
    \midrule
    \multirow{4}{*}{有效人口上界} & $N^\ast(45\,\mathrm d)$ & $4.04\times10^{4}$ & 式~\eqref{eq:dom:Nstar} \\
    & $N^\ast(60\,\mathrm d)$ & $5.37\times10^{4}$ & 式~\eqref{eq:dom:Nstar} \\
    & $N^\ast(90\,\mathrm d)$ & $8.03\times10^{4}$ & 式~\eqref{eq:dom:Nstar} \\
    & $N^\ast_\infty$ & $9.17\times10^{4}$ & 推论~\ref{cor:ceiling} \\
    \midrule
    \multirow{3}{*}{\makecell[l]{人口下界与\\附加约束}} & $N_{\rm floor}$ & $1.06\times10^{4}$ & 命题~\ref{prop:dom:floor} \\
    & $N_{\rm cum,\infty}^\ast$ & $1.18\times10^{4}$ & 推论~\ref{cor:cum} \\
    & $\max_\eta N_{\rm clr}(\eta)$ & $5166$ & 推论~\ref{cor:clr} \\
    \bottomrule
  \end{tabularx}
    \begin{tablenotes}[flushleft]
      \footnotesize
      \item 注：上段为式~\eqref{eq:dom:thetadefs} 数值求解得到的两类阈值比例；中段为命题~\ref{prop:dom:region} 的临界有效人口 $N^\ast(T_{\max})=I_{\rm peak}^{\rm T}/\theta_{\rm bind}(T_{\max})$；下段为下界与两项附加约束给出的临界值。
    \end{tablenotes}
  \end{threeparttable}
\end{table}

峰值、成本、时长和触发条件在 $(N,\eta)$ 平面上分别为
\begin{align}
  \text{峰值线：}\ &
  \eta=I_{\rm peak}^{\rm T}
  &&(\text{占优侧 }\eta\le I_{\rm peak}^{\rm T}),\\
  \text{成本线：}\ &
  \eta=\theta_{\rm cost}N
  &&(\text{占优侧 }\eta\ge\theta_{\rm cost}N),\\
  \text{时长线：}\ &
  \eta=\theta_{\rm dur}(T_{\max})N
  &&(\text{占优侧 }\eta\ge\theta_{\rm dur}(T_{\max})N),\\
  \text{触发线：}\ &
  \eta=i_{\max}^{no}N
  &&(\text{可行侧 }\eta<i_{\max}^{no}N).
\end{align}
占优区域为
\[
\mathcal{W}_{\rm pcd}
=
\bigl\{
(N,\eta):
\theta_{\rm bind}(T_{\max})N
\le\eta\le I_{\rm peak}^{\rm T},
\ \ \eta<i_{\max}^{no}N
\bigr\},
\]
比较区域的下界由成本与时长条件共同确定，其人口上限对应交点 $(N^\ast(T_{\max}),I_{\rm peak}^{\rm T})$。该点仍满足严格触发条件（$151.90$ 小于
$i_{\max}^{no}N^\ast_\infty\approx9.6\times10^{3}$）。

图~\ref{fig:dom}(a) 给出上述比较区域，图~\ref{fig:dom}(b) 加入累计感染和清零时间条件，对应集合为
\[
\mathcal{W}_{\rm cum}=\mathcal{W}_{\rm pcd}\cap\{N h\le \Itcum^{\rm T}\}\cap\{N\ge \Itcum^{\rm T}\},
\qquad
\mathcal{W}_{\rm clr}=\mathcal{W}_{\rm pcd}\cap\{t_{\rm end}\le t_{\rm end}^{\rm T}\}
\cap\{N\ge \Itcum^{\rm T}\},
\]
其中 $N\ge\Itcum^{\rm T}$ 排除了仅由人口总数就满足累计条件的情形。在本文参数下，
$\mathcal{W}_{\rm clr}\subset\mathcal{W}_{\rm cum}\subset\mathcal{W}_{\rm pcd}$。

\begin{figure}[htbp]
  \centering
  \includegraphics[width=\textwidth]{figures/fig_dom_combined.pdf}
  \caption{西安参数下 $(N_{\rm eff},\eta)$ 平面的比较区域。(a) 峰值、成本、时长和触发条件；阴影表示仅由峰值与成本限制的区域。(b) 加入累计感染和清零时间的等值曲线。竖直实线为 $N_{\rm floor}=1.06\times10^4$。累计感染边界与成本线交于 $(11762,19.48)$，清零时间等值曲线位于人口下界左侧。圆点和方点对应图~\ref{fig:panel_eta_lever}、\ref{fig:panel_N_lever} 的参数。}
  \label{fig:dom}
\end{figure}

仅考虑峰值和成本时，条件 $\theta_{\rm cost}N\le\eta\le I_{\rm peak}^{\rm T}$ 与人口下界的交集为
\begin{equation}
  N_{\rm eff}\in\bigl[\,N_{\rm floor},\,N^\ast_\infty\,\bigr]
  =\bigl[\,1.06\times10^{4},\ 9.17\times10^{4}\,\bigr],
  \label{eq:dom:main-interval}
\end{equation}
在此范围内，存在阈值使感染峰值不超过 $151.90$、加权成本不超过 $49.35$。成本边界对应的控制持续时间为 $102.83$ 天；若允许时长小于该值，有效人口上界进一步减小，见表~\ref{tab:dom:thresholds}。

在固定乘积 $\beta c_0$ 的参数分析中，区间两端随 $\beta$ 改变，但比值 $N^\ast_\infty/N_{\rm floor}$ 的相对极差仅为 $0.22\%$。因此，在所考察范围内，区间的相对宽度对 $\beta$ 的变化不敏感，具体结果见附录~\ref{app:dom:numerics}。

累计感染与清零时间条件进一步缩小了比较范围。累计感染边界与成本线的交点为 $(11762,19.48)$，相应的人口上限为 $N_{\rm cum,\infty}^\ast\approx1.18\times10^4$。与 $N_{\rm floor}$ 取交后，满足累计感染要求的区间收缩为 $N\in[1.06\times10^4,\,1.18\times10^4]$，上下端点之比不足 $1.12$。该区间接近人口下界，说明在本节设定下，累计感染条件比仅考虑峰值和成本时更严格。

清零时间边界整体位于人口下界左侧，其最大值 $N=5166$ 仍低于 $N_{\rm floor}=1.06\times10^4$。因此，在本节计算范围内，没有同时满足 $N\ge N_{\rm floor}$ 与 $t_{\rm end}\le t_{\rm end}^{\rm T}$ 的有效人口。

图~\ref{fig:panel_eta_lever} 固定 $N_{\rm eff}=2\times10^4$，比较三个阈值。各阈值均低于 TDINN 峰值参照，但累计感染均高于 TDINN。这与 $N_{\rm eff}>N_{\rm cum,\infty}^\ast$ 一致，说明峰值条件与累计感染条件并不等价。

图~\ref{fig:panel_N_lever} 固定 $\eta=100$。随 $N_{\rm eff}$ 增大，感染人数维持的阈值保持不变，持续时间增加；在本节参数范围内，$\Delta t\approx0.170N/\eta$。图~\ref{fig:panel_N_decomp} 中，人口从 $3970$ 增至 $60381$ 时，常规控制的感染峰值由 $415$ 增至 $6319$，清零时间由 $38.5$ 天增至 $50.9$ 天；阈值控制的控制持续时间由 $6.3$ 天增至 $102.8$ 天，清零时间由 $45.3$ 天增至 $201.6$ 天。较低的固定感染峰值因此可能伴随明显延长的疫情持续时间。

这些轨迹中的拐点相对位置均约为 $\lambda=0.861$，拐点处隔离率约为 $\qinf=0.4119$。前者是本组参数下的近似结果，后者由仅含 $\beta$ 的解析表达式确定。

\begin{figure}[htbp]
  \centering
  \includegraphics[width=0.8\textwidth]{figures/fig_panel_A.pdf}
  \caption{固定 $N_{\rm eff}=2\times10^4$，三个阈值 $\eta=22.7,33.1,75.2$ 分别对应 $150$ 天时长条件、成本条件和 $45$ 天时长条件。(a) 非隔离感染者人数，含同人口下的常规控制及固定 TDINN 参照；(b) 隔离率，内嵌图为累计感染。实心圆表示隔离率拐点时刻在感染人数曲线水平段上的位置，空心圆为隔离率曲线的拐点，短竖线表示清零时刻。}
  \label{fig:panel_eta_lever}
\end{figure}

\begin{figure}[htbp]
  \centering
  \includegraphics[width=0.8\textwidth]{figures/fig_panel_B.pdf}
  \caption{固定 $\eta=100$，有效人口分别为 $3970,10102,26584,60381,87944$ 时的 (a) 感染轨迹和 (b) 隔离率。灰色带表示相同人口范围内常规控制轨迹的变化范围，黑实线为固定 TDINN 参照。内嵌图中的累计感染依次约为 $945,2097,5016,10784,15410$，水平线为 TDINN 参照值 $2096.76$。圆点含义同图~\ref{fig:panel_eta_lever}。}
  \label{fig:panel_N_lever}
\end{figure}

\begin{figure}[htbp]
  \centering
  \includegraphics[width=\textwidth]{figures/fig_panel_B_trajectory_decomposition.pdf}
  \caption{固定 $\eta=100$ 时不同有效人口下的感染轨迹。(a)--(d) 分别对应 $N_{\rm eff}=3970,10102,26584,60381$，各包含阈值控制、常规控制和固定 TDINN 参照；(e) 比较四个人口规模下的常规控制；(f) 比较阈值控制与 TDINN。}
  \label{fig:panel_N_decomp}
\end{figure}

例如，$N=5\times10^4$、$\eta=100$ 时，$J=40.77$、$\Delta t=85.07$ 天，满足峰值、成本及 $T_{\max}=150$ 天的要求；但累计感染约为 $9.03\times10^3$、清零时间约为 $176$ 天，均超过 TDINN 参照值。

\subsection{常规接触率的敏感性分析}
\label{sec:dom:c0}

固定西安传播参数 $\beta,\gamma,q_0$、$N_{\rm eff}=2\times10^4$、$\theta=0.002$（即 $\eta=40$）和绝对初值 $I_0=1.00663\times10^{-3}$，考察 $c_0$ 的影响。这里的 $c_0$ 是给定的常规接触水平，不作为额外控制量。

图~\ref{fig:c0-panel} 中，各控制轨迹的感染阈值相同。增大 $c_0$ 会使启动时间提前、最大隔离率升高，与命题~\ref{prop:s1:sensitivity} 一致。$c_0=3.43$ 时，$q_{\max}=0.383<\qinf=0.4119$，无内部拐点；$c_0=3.59$ 时，$q_{\max}=0.423$，内部拐点出现。存在拐点的各条轨迹均满足 $\qinf=0.4119$。

\begin{figure}[htbp]
    \centering
    \includegraphics[width=0.8\textwidth]{c0_sensitivity_panel.pdf}
    \caption{固定 $N_{\rm eff}=2\times10^4$、$\theta=0.002$，不同 $c_0$ 下的 (a) 感染人数和 (b) 隔离率。内嵌图为累计感染。$c_0=3.43$ 时无内部拐点；其余曲线的拐点隔离率均为 $\qinf=0.4119$。圆点含义同图~\ref{fig:panel_eta_lever}。}
    \label{fig:c0-panel}
\end{figure}

图~\ref{fig:c0-scan} 中，控制持续时间随 $c_0$ 先增加后减小，在 $c_0\approx6.61$ 时达到约 $103.2$ 天；成本和累计感染的极大值分别位于 $c_0\approx18.12$ 和 $13.88$，对应 $J\approx42.80$ 和 $\Itcum\approx3610$。三个极大值并不重合。

动态清零时间在所考察范围内随 $c_0$ 增大而减小，从 $c_0=3.43$ 时约 $368$ 天降至 $c_0=30$ 时约 $102$ 天。虽然控制持续时间在部分区间增加，但启动时间和退出后的下降时间缩短，使总时间减小。因此，控制持续时间不能单独决定清零时间。

图~\ref{fig:c0-phase} 给出 $(\theta,c_0)$ 平面上的触发边界、内部拐点边界和控制持续时间驻点曲线。三者分别对应控制是否启动、隔离率是否存在内部拐点以及控制持续时间的变化方向。

\begin{figure}[htbp]
    \centering
    \includegraphics[width=0.72\textwidth]{c0_sensitivity_phase.pdf}
    \caption{$(\theta,c_0)$ 平面上的触发边界、内部拐点边界 $s^*=2\bar s$ 和控制持续时间驻点曲线 $\partial\Delta t/\partial c_0=0$。灰色等值线表示控制持续时间，水平点线为 $\theta=0.002$。}
    \label{fig:c0-phase}
\end{figure}

\subsection{初值影响与比较条件的适用范围}
\label{sec:dom:limits}

上述人口比较沿用全市数据拟合的传播参数与初值。接近 $N_{\rm floor}$ 时，$S/N$ 对 $1$ 的偏离增大，此时沿用在 $S/N\approx1$ 条件下估计的参数需要谨慎。实际干预也可能改变人群混合范围，因而阈值控制与 TDINN 未必对应相同的有效人口。

在当前数值范围内，成本、时长和累计感染边界对初值变化较不敏感，清零时间则更敏感。由 $t_1\approx r^{-1}\ln(\theta/i_0)$，其中 $r=\beta c_0(1-q_0)-\gamma=1.012\,\mathrm{d}^{-1}$，$i_0$ 每变化十倍，启动时间约变化 $2.28$ 天。若改用固定归一化初值 $i_0=7.65\times10^{-11}$，清零边界对应的最大人口由 $5166$ 变为 $2439$。两者均低于 $N_{\rm floor}=1.06\times10^4$，故在这两种设定下，加入清零条件后的比较区域均为空。

\section{讨论与局限}
\label{sec:discussion}

本文在含接触追踪与隔离的 SIQR 型模型中，研究了接触率保持不变的隔离率阈值控制。由常规阶段的首次积分和恒定感染条件，得到启动点、阈值控制期轨迹、隔离控制函数及控制持续时间，并给出动态清零时间和累计感染的分段表达。控制在启动时使隔离率跳跃增加，随后逐渐放松，直至恢复常规水平。

容量阈值与控制时间之间存在直接的权衡。降低 $\eta$ 可限制感染峰值，但会延长阈值控制期并增加隔离成本。西安参数下，$\eta=0.002N$ 时阈值控制的加权成本低于 TDINN，感染峰值、累计感染和清零时间却更高。因此，策略比较需要同时考察这些指标，而不能仅依据成本或峰值作出判断。

在固定归一化初值和其余参数时，控制持续时间、成本和最大隔离率仅通过 $\theta=\eta/N$ 依赖于人口与阈值。固定绝对初值的数值结果在所考察范围内近似满足这一关系。累计感染和清零时间还涉及绝对人口规模，这也是加入相应条件后比较区域明显缩小的原因。

有效人口分析以 TDINN 的全市拟合结果为固定参照，所得区间依赖于相同混合人口假设和当前成本定义。边界明确的校园、厂区或邮轮，与通过空间分区和接触限制形成的人群，其成本含义并不相同；后一情形的分区成本未计入 $J$。因而 $J_c=0$ 只表示阈值策略未进一步降低模型中的接触率，不表示实际干预没有其他成本。有效人口仍需由接触调查、出行数据或实际防控单元等外部信息确定。

本文采用由理论轨迹预先计算的时间控制函数，尚未考虑观测噪声和参数不确定性。西安阈值控制轨迹属于给定模型下的反事实计算，不是经观测验证的实际预测。后续可在其他地区或病种的数据上检验相关结果，并考察联合调节接触率与隔离率时的变化。

\clearpage
\appendix
\section{不同有效人口下的阈值比较}
\label{app:panel_eta_N}

本附录采用三个阈值比例：
\[
\theta_{\rm dur150}=1.137\times10^{-3},\qquad
\theta_{\rm cost}=1.656\times10^{-3},\qquad
\theta_{\rm dur45}=3.762\times10^{-3}.
\]
初值仍取固定绝对值 $I_0$，不同 $N_{\rm eff}$ 下不重新拟合。三个人口规模分别对应累计感染、$45$ 天控制持续时间和成本条件的临界值，结果见图~\ref{fig:panel_eta_Ncumstar}--\ref{fig:panel_eta_Nstarinf}。各参数组中，拐点的相对位置约为 $\lambda=0.861$，而绝对时刻随启动时间改变。

\begin{figure}[!htbp]
  \centering
  \includegraphics[width=0.7\textwidth]{figures/fig_panel_A_Ncumstar.pdf}
  \caption{$N_{\rm eff}=11763\approx N^\ast_{\rm cum,\infty}$ 时的阈值比较。$\eta=13.375,19.481,44.246$，对应清零时间约为 $227.5,175.3,107.1$ 天。成本条件对应的累计感染约为 TDINN 参照值 $2096.76$；另外两个阈值下的累计感染分别低于和高于该参照。}
  \label{fig:panel_eta_Ncumstar}
\end{figure}

\begin{figure}[!htbp]
  \centering
  \includegraphics[width=0.7\textwidth]{figures/fig_panel_A_Nstar45.pdf}
  \caption{$N_{\rm eff}=40382\approx N^\ast(45\,\mathrm d)$ 时的阈值比较。$45$ 天时长条件对应 $\eta=151.904$，约等于 TDINN 峰值；其余阈值 $45.918$ 和 $66.880$ 均低于该参照。}
  \label{fig:panel_eta_Nstar45}
\end{figure}

\begin{figure}[!htbp]
  \centering
  \includegraphics[width=0.7\textwidth]{figures/fig_panel_A_Nstarinf.pdf}
  \caption{$N_{\rm eff}=91721\approx N^\ast_\infty$ 时的阈值比较。成本条件对应 $\eta=151.904$，约等于 TDINN 峰值；$45$ 天时长条件对应 $\eta=345.02$，超过该峰值；$150$ 天时长条件对应 $\eta=104.295$，低于该峰值。}
  \label{fig:panel_eta_Nstarinf}
\end{figure}

\clearpage

\section{传播概率的影响与比较边界的数值结果}
\label{app:dom:numerics}

本附录给出传播概率变化及累计感染、清零时间比较边界的补充结果。

固定乘积 $\beta c_0=1.9305$，在 $\beta\in[0.10,0.30]$ 内，人口区间两端随 $\beta$ 改变，但端点比值变化较小：
$N^\ast_\infty/N_{\rm floor}\in[8.669,\,8.688]$，极差 $0.22\%$（$\beta=0.10/0.1498/0.20/0.30$ 处
$N^\ast_\infty=13.48/9.17/7.00/4.83\times10^4$、$N_{\rm floor}=1.55/1.06/0.81/0.56\times10^4$）。
这一结果并非精确比例关系：
$\beta N^\ast_\infty$ 与 $\beta N_{\rm floor}$ 各自随 $\beta$ 漂移约 $7.6\%$ 与 $7.8\%$，
二者的变化幅度接近，但都不严格正比于 $1/\beta$。

在 $\eta\in[10,151.90]$ 内，累计感染等值曲线 $N h(\theta,N)=\Itcum^{\rm T}$ 从 $(9471,151.90)$ 延伸至 $(12177,10)$，$\eta=100$ 时 $N_{\rm cum}\approx1.010\times10^4$。它与成本线的交点为 $(11762,19.48)$，给出 $N_{\rm cum,\infty}^\ast\approx1.18\times10^4$，约为 $N^\ast_\infty$ 的 $0.13$。现有 $16$ 个计算点及其插值仅显示一次交点，尚不能据此确定一般参数下交点的唯一性。

$t_{\rm end}=t_{\rm end}^{\rm T}$ 在同一阈值范围内从 $(865.9,\,10)$ 延伸到 $(5165.7,\,151.90)$，
$\eta=100$ 处 $N_{\rm clr}\approx3.97\times10^{3}$，全段最大值为 $N=5166$。

第~\ref{sec:scaling}~小节的累计比例计算至 $t_2$，本附录中的 $h$ 则计算至动态清零时刻，包含退出后的感染过程。两者的积分终点不同。

\clearpage

\section{常规接触率的敏感性分析}
\label{app:c0_scan}

固定 $N_{\rm eff}=2\times10^4$、$\theta=0.002$，图~\ref{fig:c0-scan} 给出各项指标随 $c_0$ 的变化。控制持续时间、成本与累计感染均存在数值极大值，其余时间指标在所考察范围内递减。

\begin{figure}[htbp]
    \centering
    \includegraphics[width=\textwidth]{c0_sensitivity_scan.pdf}
    \caption{固定 $N_{\rm eff}=2\times10^4$、$\theta=0.002$ 时各指标随 $c_0$ 的变化。(a)--(d) 依次为启动时间、控制持续时间、退出后的下降时间和动态清零时间；(e)--(h) 依次为最大隔离率、拐点相对位置、成本和累计感染。竖线标示触发边界 $3.334$、拐点边界 $3.543$ 和时长极大值位置 $6.607$。}
    \label{fig:c0-scan}
\end{figure}


\section{接触率与传播概率对拐点存在性的影响}
\label{app:c0_beta}

固定 $\theta=0.002$、$q_0=0.3230$ 及绝对初值，图~\ref{fig:c0-beta-existence} 给出 $(c_0,\beta)$ 平面上满足 $s_c<2\bar s<s^*$ 的参数范围。西安参数 $(12.8872,0.1498)$ 位于内部拐点存在区域内。

\begin{figure}[htbp]
    \centering
    \includegraphics[width=0.7\textwidth]{c0_beta_existence.pdf}
    \caption{固定 $\theta=0.002$、$q_0=0.3230$ 时，$(c_0,\beta)$ 平面上的内部拐点存在范围。边界为 $2\bar s=s^*$ 和 $\beta_{\max}=0.2614$，标记点为西安参数 $(12.8872,0.1498)$。}
    \label{fig:c0-beta-existence}
\end{figure}

\clearpage

% 参考文献
\printbibliography[title={参考文献}]


\end{document}

